%  Informe de Practica Profesional I - Implementacion de nuevos
%  algoritmos de agrupamiento en BHAVRP
%  Compilar con: pdflatex (dos veces para el indice)
%==============================================================
\documentclass[12pt,letterpaper,oneside]{report}

%---------------------- Codificacion e idioma -------------------
\usepackage[utf8]{inputenc}
\usepackage[T1]{fontenc}
\usepackage[spanish,es-noshorthands,es-tabla]{babel}

%---------------------- Fuentes (Times) -------------------------
\usepackage{newtxtext}
\usepackage{amsmath,amssymb}
\usepackage{newtxmath}

%---------------------- Pagina y margenes -----------------------
\usepackage[letterpaper,top=2.5cm,bottom=2.5cm,left=3cm,right=2.5cm]{geometry}
\usepackage{setspace}
\onehalfspacing                 % interlineado 1.5
\usepackage{indentfirst}        % sangria en primer parrafo
\setlength{\parindent}{1.25cm}
\setlength{\parskip}{0pt}

%---------------------- Graficos y figuras ----------------------
\usepackage{graphicx}
\usepackage{float}
\graphicspath{{Figuras/}}
% Las figuras del documento original deben exportarse del Word
% (clic derecho -> Guardar como imagen) dentro de la carpeta Figuras/.

%---------------------- Tablas ----------------------------------
\usepackage{booktabs}
\usepackage{array}
\usepackage{multirow}
\usepackage{colortbl}
\usepackage{tabularx}

%---------------------- Encabezados -----------------------------
\usepackage{fancyhdr}
\pagestyle{fancy}
\fancyhf{}
\fancyhead[R]{\thepage}
\renewcommand{\headrulewidth}{0pt}

%---------------------- Titulos de capitulos/secciones ----------
\usepackage{titlesec}
\titleformat{\chapter}[block]
  {\normalfont\Large\bfseries}{Cap\'itulo \thechapter:}{0.6em}{}
\titlespacing*{\chapter}{0pt}{-20pt}{20pt}
\titleformat{\section}{\normalfont\large\bfseries}{\thesection}{0.8em}{}
\titleformat{\subsection}{\normalfont\normalsize\bfseries}{\thesubsection}{0.8em}{}

%---------------------- Captions --------------------------------
\usepackage{caption}
\captionsetup[table]{position=above,skip=6pt}
\captionsetup[figure]{position=below,skip=10pt}
\captionsetup{font=small,labelfont=bf}

%---------------------- Hipervinculos ---------------------------
\usepackage[plainpages=false,pdfpagelabels ,hidelinks]{hyperref}

%---------------------- Comando util ----------------------------
% Inserta la figura si existe el archivo; si no, deja un recuadro
% indicativo con el nombre del fichero que falta.
\newcommand{\figura}[4]{%
  \begin{figure}[H]
    \centering
    \IfFileExists{Figuras/#1}{%
      \includegraphics[width=#2]{#1}%
    }{%
      \fbox{\parbox[c][6cm][c]{0.8\textwidth}{\centering
        \textit{[Insertar aqu\'i la figura: \texttt{ \detokenize{Figuras/#1}}]}}}}%
    \caption{#3}\label{#4}
  \end{figure}%
}

\begin{document}
\pagenumbering{roman}

%==============================================================
%                         PORTADA
%==============================================================
\begin{titlepage}
\centering
{\large \textbf{Universidad Tecnol\'ogica de La Habana}}\\[2pt]
{\large \textbf{``Jos\'e Antonio Echeverr\'ia''}}\\[2pt]
{\large \textbf{Facultad de Ingenier\'ia Inform\'atica}}\\[40pt]

\IfFileExists{Figuras/isotipo.png}{%
  \includegraphics[width=3.2cm]{isotipo.png}\\[40pt]%
}{%
  \fbox{\parbox[c][3.2cm][c]{3.2cm}{\centering\textit{[Logo: Figuras/isotipo.png]}}}\\[40pt]%
}

{\Large \textbf{Informe de Pr\'actica Profesional I}}\\[30pt]

{\large \textbf{Implementaci\'on de nuevos algoritmos de agrupamiento en BHAVRP}}\\[45pt]

\begin{flushleft}
\hspace{2cm}{\large \textbf{Autor:}}\\[6pt]
\hspace{2cm}{\large Maykel Kevin D\'iaz Men\'endez}\\[25pt]
\hspace{2cm}{\large \textbf{Tutores:}}\\[6pt]
\hspace{2cm}{\large Ing.\ Eric Ramos Arag\'on}\\
\hspace{2cm}{\large Ing.\ Ananda de la Caridad Morales Morales}
\end{flushleft}

\vfill
{\large \textbf{La Habana, Cuba}}\\[6pt]
{\large \textbf{Marzo 2026}}
\end{titlepage}

%==============================================================
%                         RESUMEN
%==============================================================
\chapter*{Resumen}

El Problema de Planificaci\'on de Rutas de Veh\'iculos con M\'ultiples Dep\'ositos (MDVRP) constituye una variante compleja dentro de los problemas de optimizaci\'on combinatoria. En la pr\'actica, este problema suele resolverse dividi\'endolo en dos etapas: asignar los clientes a los dep\'ositos y construir las rutas desde cada dep\'osito. La Biblioteca de Heur\'isticas de Asignaci\'on para Problemas de Planificaci\'on de Rutas de Veh\'iculos con M\'ultiples Dep\'ositos (BHAVRP) es un componente de software dise\~nado para resolver esta primera fase de asignaci\'on. La biblioteca presenta una arquitectura modular que encapsula m\'etodos de agrupamiento clasificados en particionales y jer\'arquicos, incluye un sistema de tratamiento de excepciones y un conjunto de m\'etricas de evaluaci\'on internas. Actualmente, BHAVRP carece de algoritmos de agrupamiento basados en densidad y cuenta con un solo algoritmo de agrupamiento jer\'arquico. Se propone la incorporaci\'on de dos nuevos algoritmos de agrupamiento: DBSCAN como basado en densidad, y CURE como jer\'arquico. Ambos algoritmos de agrupamiento son adaptados para resolver la fase de asignaci\'on de clientes a dep\'ositos incorporando restricciones de capacidad de los dep\'ositos y mecanismos de asignaci\'on controlada. Por \'ultimo, se realiza un estudio experimental que permite garantizar el correcto funcionamiento de los m\'odulos implementados, as\'i como identificar qu\'e algoritmos generan asignaciones de mayor calidad. Para ello, se utilizan los criterios definidos por las m\'etricas de evaluaci\'on que ya tiene BHAVRP.

\vspace{12pt}
\noindent\textbf{Palabras clave:} algoritmos de agrupamiento, BHAVRP, CURE, DBSCAN, MDVRP.

%==============================================================
%                         ABSTRACT
%==============================================================
\chapter*{Abstract}

The Multi-Depot Vehicle Routing Problem (MDVRP) is a complex variant within combinatorial optimization problems. In practice, this problem is often solved by dividing it into two stages: assigning customers to depots and building routes from each depot. The Assignment Heuristics Library for Multi-Depot Vehicle Routing Problems (BHAVRP) is a software component designed to solve this first assignment phase. The library features a modular architecture that encapsulates clustering methods classified into partitional and hierarchical, includes an exception handling system, and a set of internal evaluation metrics. Currently, BHAVRP lacks density-based clustering algorithms and has only one hierarchical clustering algorithm. The incorporation of two new clustering algorithms is proposed: DBSCAN as a density-based algorithm, and CURE as a hierarchical algorithm. Both clustering algorithms are adapted to solve the customer-to-depot assignment phase by incorporating depot capacity constraints and controlled assignment mechanisms. Finally, an experimental study is conducted to ensure the correct functioning of the implemented modules, as well as to identify which algorithms generate higher-quality assignments. For this purpose, the criteria defined by the evaluation metrics already present in BHAVRP are used.

\vspace{12pt}
\noindent\textbf{Keywords:} BHAVRP, CURE, clustering algorithms, DBSCAN, MDVRP.

%==============================================================
%                       INDICES
%==============================================================
\tableofcontents
\listoffigures
\listoftables

\clearpage
\pagenumbering{arabic}
\setcounter{page}{1}

%==============================================================
%                      INTRODUCCION
%==============================================================
\chapter*{Introducci\'on}
\addcontentsline{toc}{chapter}{Introducci\'on}
\markboth{Introducci\'on}{}

La gesti\'on de la distribuci\'on, el transporte y la planificaci\'on de rutas son temas de gran importancia que han sido abordados a nivel internacional. Son ramas de la log\'istica donde yacen numerosas oportunidades de mejora y por esta raz\'on ha recibido especial atenci\'on motivando un intenso trabajo te\'orico en el desarrollo de modelos y algoritmos. En este contexto surgen los Problemas de Planificaci\'on de Rutas de Veh\'iculos (\textit{Vehicle Routing Problems}, VRP) [1].

En [2] se introduce el \textit{``Truck Dispatching Problem''}, modelando c\'omo una flota de camiones homog\'eneos podr\'ia satisfacer la demanda de petr\'oleo de varias estaciones de servicio desde un dep\'osito central y con una distancia m\'inima recorrida. Cinco a\~nos despu\'es, en [3] se generaliza este a un problema de optimizaci\'on lineal que se encuentra com\'unmente en el \'ambito de la log\'istica y el transporte. El VRP es un problema NP-duro [4] por lo que se consideran muy dif\'iciles de resolver. Los problemas de la vida real son considerablemente m\'as grandes y complejos que los escenarios acad\'emicos, por lo que las heur\'isticas y metaheur\'isticas son necesarias para obtener buenos resultados. Estos algoritmos no garantizan encontrar la soluci\'on \'optima, pero s\'i proporcionan tiempos computacionales razonables. La cantidad de m\'etodos de soluci\'on introducidos en la literatura acad\'emica ha crecido en las \'ultimas d\'ecadas [5].

El Problema de Planificaci\'on de Rutas de Veh\'iculos con M\'ultiples Dep\'ositos (\textit{Multi-Depot Vehicle Routing Problems}, MDVRP) es una variante del VRP cl\'asico que incorpora varios dep\'ositos con una ubicaci\'on predefinida. Cada dep\'osito cuenta con una flota de veh\'iculos de capacidad homog\'enea que tiene como objetivo satisfacer la demanda de un conjunto de clientes dispersos geogr\'aficamente. Cada veh\'iculo debe comenzar y terminar su ruta en el mismo dep\'osito, y cada cliente debe ser visitado una \'unica vez. Teniendo en cuenta la distancia entre la ubicaci\'on de los clientes y dep\'ositos, el MDVRP busca asignar un cliente a cada dep\'osito y se construye una ruta para cada veh\'iculo dentro de un dep\'osito para atender dichos clientes. El objetivo al crear la ruta es reducir al m\'inimo el costo total del recorrido [6].

Durante un largo per\'iodo de tiempo, los m\'etodos heur\'isticos basados en la b\'usqueda tab\'u han sido los m\'as empleados en la literatura [6]. Sin embargo, este problema genera una gran cantidad de combinaciones que se refleja en los altos tiempos computacionales para obtener la soluci\'on. El potencial de las t\'ecnicas de an\'alisis de cl\'usteres para la soluci\'on de los MDVRP ha sido reconocido por autores como [7], [2], [8] y [9]. El an\'alisis de cl\'usteres estudia la divisi\'on de objetos en grupos basado en una o varias caracter\'isticas. Un grupo se define como una zona con muchos objetos parecidos muy cercanos entre s\'i, que est\'a separada de otros grupos por zonas donde casi no hay objetos. Esta definici\'on de grupo es bastante acorde a la agrupaci\'on de clientes por lo que se propone usar el an\'alisis de cl\'usteres en la soluci\'on del MDVRP, espec\'ificamente en la fase de asignaci\'on de clientes a dep\'ositos [7]. En [10] se describe la t\'ecnica ``divide y vencer\'as'', utilizada para reducir la complejidad computacional en el empleo heur\'isticas para resolver el MDVRP.

Sin embargo, a\'un existen desaf\'ios en cuanto a la eficiencia computacional debido a los altos tiempos y la calidad de las agrupaciones dependiendo del escenario [7]. En la literatura sobresalen algoritmos como \textit{Density-Based Spatial Clustering of Applications with Noise} (DBSCAN) que permiten identificar cl\'usteres basados en la densidad local [11]. Por otra parte, \textit{Clustering Using Representatives} (CURE) identifica variaciones en la forma de los cl\'usteres mediante el uso de puntos representativos [12].

En este contexto, la Universidad Tecnol\'ogica de la Habana ``Jos\'e Antonio Echeverr\'ia'' (CUJAE) ha desarrollado la Biblioteca de Heur\'isticas de Asignaci\'on para Problemas de Planificaci\'on de Rutas de Veh\'iculos con M\'ultiples Dep\'ositos (BHAVRP). Este componente est\'a dise\~nado para resolver la fase de asignaci\'on en MDVRP utilizando una arquitectura modular y 20 algoritmos repartidos en heur\'isticas cl\'asicas de asignaci\'on y algoritmos de agrupamiento de Miner\'ia de Datos adaptados al contexto de los MDVRP [13]. Esta alternativa gratuita y de c\'odigo abierto supone una ventaja frente a otras herramientas que no tienen una colecci\'on de algoritmos tan extensa o son software comercial.

Esta biblioteca, en la categor\'ia de algoritmos de agrupamiento, incluye algoritmos particionales, algoritmos jer\'arquicos y un algoritmo basado en densidad. En la literatura se ha sugerido que ciertas m\'etricas de validaci\'on internas como Dunn Index o Silhouette podr\'ian tener una correlaci\'on con la calidad de las rutas en la posterior fase de planificaci\'on. En investigaciones anteriores [14], [15], se realiz\'o un an\'alisis experimental donde se concluy\'o que no hab\'ia una correlaci\'on positiva significativa entre las m\'etricas internas dentro de BHAVRP y la calidad de las rutas obtenidas. Sin embargo, este estudio fue realizado con poca variedad de algoritmos de agrupamiento y sin considerar una amplia variedad de heur\'isticas de construcci\'on.

Debido a la situaci\'on problem\'atica antes descrita, se identifica como problema a resolver la falta de un estudio exhaustivo que determine si existe alguna m\'etrica interna que correlacione positivamente con la calidad de las rutas obtenidas en la planificaci\'on, lo que impide establecer un criterio de selecci\'on de los algoritmos de BHAVRP.

Para dar soluci\'on al problema planteado se traza como objetivo general: Desarrollar una nueva versi\'on de BHAVRP que incluya m\'as algoritmos de agrupamiento basados en densidad.

Para dar cumplimiento al objetivo general se definieron los siguientes objetivos espec\'ificos con sus respectivas tareas:

\begin{enumerate}
  \item Identificar algoritmos de agrupamiento utilizados para resolver la fase de asignaci\'on de los problemas de planificaci\'on de rutas de veh\'iculos con m\'ultiples dep\'ositos.
  \begin{itemize}
    \item Realizar b\'usqueda bibliogr\'afica de las caracter\'isticas de los problemas de planificaci\'on de rutas de veh\'iculos con m\'ultiples dep\'ositos y sus aplicaciones pr\'acticas.
    \item Realizar b\'usqueda bibliogr\'afica de las propuestas de soluci\'on de MDVRP donde se realice la fase de asignaci\'on de clientes a dep\'ositos usando como v\'ia de resoluci\'on los algoritmos de agrupamiento.
  \end{itemize}
  \item Implementar en BHAVRP un nuevo algoritmo de agrupamiento.
  \begin{itemize}
    \item Analizar algoritmos de agrupamiento en cuanto a su funcionamiento para su incorporaci\'on en la nueva versi\'on de la biblioteca.
    \item Seleccionar los algoritmos de agrupamiento que se van a incorporar a la biblioteca para resolver la variante MDVRP.
    \item Implementar y adaptar el algoritmo seleccionado teniendo en cuenta las caracter\'isticas que presentan.
  \end{itemize}
  \item Validar el funcionamiento de estos algoritmos en BHAVRP mediante la ejecuci\'on de experimentos.
  \begin{itemize}
    \item Realizar pruebas unitarias, donde se muestre que el algoritmo implementado funcione correctamente.
    \item Comparar los resultados obtenidos por los algoritmos de agrupamiento propuestos con los algoritmos de agrupamiento implementados en la versi\'on anterior de la biblioteca.
  \end{itemize}
\end{enumerate}

El cumplimiento del objetivo general tiene como valor pr\'actico un aporte significativo en la resoluci\'on de la fase de asignaci\'on del problema de planificaci\'on de rutas de veh\'iculos con m\'ultiples dep\'ositos. Los resultados tienen valor cient\'ifico y tecnol\'ogico ya que se pone a disposici\'on de la comunidad cient\'ifica un componente de software que brinda nuevos algoritmos para la soluci\'on de MDVRP.

Para alcanzar los objetivos mencionados, el trabajo se encuentra estructurado en tres cap\'itulos, incluyendo secciones importantes como conclusiones, recomendaciones y referencias bibliogr\'aficas. El Cap\'itulo 1: Fundamentos te\'oricos del problema de planificaci\'on de rutas y su soluci\'on mediante diferentes m\'etodos donde se presenta el estado del arte del MDVRP, las estrategias de soluci\'on existentes, as\'i como un an\'alisis comparativo de bibliotecas. El Cap\'itulo 2: Incorporaci\'on de algoritmos a la Biblioteca de Heur\'isticas de Asignaci\'on del Problema de Planificaci\'on de Rutas con M\'ultiples Dep\'ositos describe los algoritmos de agrupamiento de BHAVRP, los algoritmos propuestos detallando el funcionamiento de estos seg\'un la literatura y las adaptaciones al MDVRP que hay que realizar. Finalmente, el Cap\'itulo 3: An\'alisis experimental expone varias pruebas unitarias de las funciones m\'as relevantes de los algoritmos implementados, as\'i como un estudio comparativo entre algoritmos de asignaci\'on aplicados sobre ciertas instancias y pruebas estad\'isticas no param\'etricas para determinar su rendimiento y aplicabilidad.
%==============================================================
%                    CAPITULO 1
%==============================================================
\chapter{Fundamentos te\'oricos del problema de planificaci\'on de rutas y su soluci\'on mediante diferentes m\'etodos}

En este cap\'itulo se describen los aspectos te\'oricos de la investigaci\'on que se encuentran relacionados con los MDVRP. Entre los aspectos que se abordan est\'an las caracter\'isticas y variantes del problema, as\'i como los diferentes m\'etodos de soluci\'on encontrados en la literatura. Adem\'as, se hace \'enfasis en los algoritmos de agrupamiento como alternativa para la fase de asignaci\'on de clientes a dep\'ositos. Por \'ultimo, se describe la Biblioteca de Heur\'isticas de Asignaci\'on para Problemas de Planificaci\'on de Rutas de Veh\'iculos con M\'ultiples Dep\'ositos y sus algoritmos existentes.

\section{Problemas de Planificaci\'on de Rutas de Veh\'iculos}

El concepto de Problema de Planificaci\'on de Rutas de Veh\'iculos (\textit{Vehicle Routing Problem}, VRP) es una generalizaci\'on del Problema del Viajero (\textit{Traveling Salesman Problem}, TSP). En concreto, el VRP es un problema de optimizaci\'on combinatoria cuyo objetivo es atender las demandas de clientes que se encuentran dispersos geogr\'aficamente sujetos a un conjunto de restricciones. La diferencia radica en que TSP dispone de un mapa con m\'ultiples puntos que un individuo tiene que visitar para volver a su punto de origen realizando la ruta m\'as corta posible. En cambio, el VRP no se limita a solo una persona, sino que pueden existir m\'as de un individuo, en este caso veh\'iculos, que normalmente deben entregar objetos en diversos puntos de la ciudad a clientes. De este modo, el problema se vuelve m\'as complejo que el TSP [16].

Las aplicaciones inform\'aticas para el dise\~no de rutas se derivan de modelos VRP que atienden diferentes necesidades e incluyen diversas restricciones. Esto ha dado como resultado el surgimiento de diferentes variantes del cl\'asico VRP, de las cuales se pueden mencionar las siguientes:

\begin{itemize}
  \item VRP con recolecci\'on y entrega (VRPPD): estudia el caso donde una empresa debe recolectar y entregar bienes en cantidades espec\'ificas para cada cliente visitado [17].
  \item VRP con flota heterog\'enea: los veh\'iculos tienen diferentes capacidades de carga [18].
  \item \textit{Open VRP} (OVRP): se refiere al caso de cuando el veh\'iculo no regresa al dep\'osito de origen, es decir, que se dirige a otro dep\'osito o sitio diferente para cargar productos, en un contexto de flota rentada [19].
  \item VRP con ventanas de tiempo (VRPTW): plantea que cada cliente tiene que ser atendido de manera obligada dentro de un cierto horario espec\'ifico [20].
  \item VRP con dep\'ositos m\'ultiples (MDVRP): implica que la empresa posee diversos dep\'ositos desde los cuales puede abastecer a los clientes [21].
  \item \textit{Stochastic VRP} (SVRP): Propone que las variables involucradas son desconocidas o aleatorias, e incluso que se pueden encontrar en un rango de probabilidad, tales como el n\'umero de clientes, sus demandas, tiempo de servicio, tiempo de viaje, entre otras [22].
  \item \textit{Periodic VRP} (PVRP): analiza un problema donde los pedidos pueden ser llevados s\'olo en ciertos d\'ias [23].
  \item Capacitated VRP: estudia el caso donde los veh\'iculos tienen una capacidad m\'axima de carga y se deben dise\~nar rutas que satisfagan la demanda de los clientes sin exceder dicha capacidad [24].
\end{itemize}

\subsection{Problema de Planificaci\'on de Rutas de Veh\'iculos con M\'ultiples Dep\'ositos}

El MDVRP es un problema de optimizaci\'on combinatoria, categorizado por su complejidad como NP-Duro [25]. Por lo que es dif\'icil encontrar su soluci\'on \'optima usando m\'etodos exactos en un tiempo polinomial. Esta variante cuenta con varios dep\'ositos y cada uno tiene su propia flota de veh\'iculos [26].

El problema puede ser representado mediante modelos de optimizaci\'on. Sea el grafo completo $G = (V,A)$, donde el conjunto de v\'ertices $V = \{v_{01},...,v_{0d}\} \cup \{v_{1},...,v_{n}\}$ incluye los $d$ dep\'ositos y $c$ clientes y el conjunto de arcos $A = \{(i,j) \vee i,j \in V, i \neq j\}$ [27]. El objetivo es minimizar el costo total de las rutas, sujeto a las restricciones operativas. La funci\'on objetivo, variables de decisi\'on y las restricciones fundamentales son [28]:

\noindent\textbf{Funci\'on Objetivo:} Minimizar el costo total de transporte, generalmente expresado como la suma de las distancias recorridas por todos los veh\'iculos.

\begin{equation}
\min \sum_{k \in K}{\sum_{(i,j) \in A}{c_{ij} \cdot x_{ijk}}}
\end{equation}

\noindent\textbf{Variables de Decisi\'on:}
\begin{equation}
x_{ijk} = \left\{ \begin{array}{ll}
1 & \text{si el veh\'iculo } k \text{ viaja del v\'ertice } i \text{ al v\'ertice } j \\
0 & \text{en caso contrario}
\end{array} \right. \quad \forall(i,j) \in A,\ \forall k \in K
\end{equation}

\noindent\textbf{Restricciones:}

Cada ruta debe comenzar y terminar en el mismo dep\'osito.
\begin{equation}
\text{Salida del Dep\'osito:}\quad \sum_{j \in C}{x_{djk}} = 1 \quad \forall d \in V_{0},\ \forall k \in K
\end{equation}
\begin{equation}
\text{Entrada del Dep\'osito:}\quad \sum_{j \in C}{x_{jdk}} = 1 \quad \forall d \in V_{0},\ \forall k \in K
\end{equation}

\begin{itemize}
  \item Cada cliente debe ser visitado exactamente una vez.
\end{itemize}
\begin{equation}
\sum_{k \in K}{\sum_{j \in V, j \neq i}{x_{ijk}}} = 1 \quad \forall i \in C
\end{equation}

\begin{itemize}
  \item La demanda del cliente es un n\'umero diferente a cero y no negativo.
\end{itemize}
\begin{equation}
q_{i} > 0 \quad \forall i \in C
\end{equation}

\begin{itemize}
  \item La demanda total de cada ruta no debe exceder la capacidad del veh\'iculo.
\end{itemize}
\begin{equation}
\sum_{i \in C}{q_{i} \cdot \left( \sum_{j \in V, j \neq i}{x_{ijk}} \right)} \leq Q \quad \forall k \in K
\end{equation}

\begin{itemize}
  \item Todos los veh\'iculos tienen una misma capacidad.
\end{itemize}
\begin{equation}
Q_{k} = Q \quad \forall k \in K
\end{equation}

\begin{itemize}
  \item Se dispone de un n\'umero limitado de veh\'iculos.
\end{itemize}
\begin{equation}
\sum_{k \in K}{\sum_{j \in C}{x_{djk}}} \leq m_{d} \quad \forall d \in V_{0}
\end{equation}

\begin{itemize}
  \item Todos los dep\'ositos tienen el mismo n\'umero de veh\'iculos disponibles.
\end{itemize}
\begin{equation}
m_{d} = m \quad \forall d \in V_{0} \qquad \sum_{k \in K}{\sum_{j \in C}{x_{djk}}} \leq m \quad \forall d \in V_{0}
\end{equation}

En la Figura~\ref{fig:mdvrp} se ilustra la estructura del problema:

\begin{itemize}
  \item Se observan tres dep\'ositos (representados por cuadrados), cada uno con una flota de veh\'iculos.
  \item Los clientes (c\'irculos rojos) son asignados a un dep\'osito espec\'ifico.
  \item Las l\'ineas verdes muestran las rutas de los veh\'iculos, que parten y regresan a su dep\'osito de origen.
\end{itemize}

\figura{fig01_mdvrp.png}{0.75\textwidth}{Representaci\'on gr\'afica de MDVRP.}{fig:mdvrp}

La literatura documenta m\'ultiples aplicaciones del MDVRP a problemas log\'isticos reales. Un estudio aplicado a una red de distribuci\'on de productos de limpieza y cuidado personal en Turqu\'ia utiliza el MDVRP para optimizar las rutas desde tres dep\'ositos a 502 clientes. Mediante un enfoque de agrupar primero y rutear despu\'es se logran mejoras significativas en la eficiencia de la distribuci\'on, utilizando datos reales de demanda y coordenadas geogr\'aficas [29].

Otro ejemplo se encuentra en el Reino Unido, donde se desarrolla un modelo de programaci\'on entera mixta (MIP) [30] para resolver un MDVRP que incluye complejidades del mundo real como una flota heterog\'enea, costos de conductores y ventanas de tiempo. La implementaci\'on de este modelo permite reducir sustancialmente los costos de transporte sin necesidad de abrir nuevas instalaciones [31].

Tambi\'en se encuentra un caso de estudio en Estambul que trata sobre la distribuci\'on de una empresa de bebidas utilizando datos hist\'oricos de GPS para construir una matriz de tiempo-distancia m\'as precisa. Se eval\'uan m\'ultiples enfoques de optimizaci\'on, incluyendo algoritmos evolutivos y de aprendizaje profundo, para resolver el MDVRP. Los resultados muestran que estos m\'etodos son efectivos para manejar la complejidad y la naturaleza din\'amica de las entregas en entornos urbanos de alta densidad [32].

\section{Algoritmos heur\'isticos}

Debido a la complejidad que presentan los VRP, existen dos tipos de algoritmos que pueden resolver este tipo de problemas, los m\'etodos exactos y aproximados. Cada uno de ellos posee diferentes caracter\'isticas que pueden considerarse ventajas y/o desventajas que deben ser tomadas en cuenta al ser utilizados [26]. En la Figura~\ref{fig:tecnicas} se puede apreciar una taxonom\'ia de estos m\'etodos:

\figura{fig02_tecnicas.png}{0.85\textwidth}{T\'ecnicas utilizadas para la soluci\'on de MDVRP.}{fig:tecnicas}

Entre los m\'etodos exactos se pueden distinguir tres grupos: b\'usqueda directa de \'arbol, programaci\'on din\'amica y programaci\'on lineal, entera y mixta. Todos estos m\'etodos son eficientes en problemas de hasta 50 clientes, debido a restricciones de tiempo computacional. Se han llevado a cabo algunas investigaciones introduciendo m\'etodos exactos para MDVRP [15], [26]. Algunos ejemplos de los m\'as utilizados para el MDVRP de este tipo de m\'etodos son:

\begin{itemize}
  \item Ramificaci\'on y Corte: los algoritmos \textit{Branch \& Cut} se basan en el c\'alculo simult\'aneo de cotas inferiores y superiores de la soluci\'on \'optima. Las cotas superiores se obtienen por medio de la resoluci\'on de las relajaciones lineales asociadas a cada nodo del \'arbol de enumeraci\'on, eventualmente ajustados por medio de planos de corte, mientras que las cotas inferiores son halladas por medio de procedimientos heur\'isticos y cuando la relajaci\'on lineal asociada con alg\'un subproblema tiene \'optimo entero [33].
  \item Programaci\'on Lineal: la programaci\'on lineal es un m\'etodo a trav\'es del cual se optimiza una funci\'on objetivo, bien sea maximizando o minimizando dicha funci\'on, en la cual las variables est\'an elevadas a la potencia 1. Este campo de la programaci\'on lineal es un \'area de la programaci\'on matem\'atica, dedicada a maximizar o minimizar una funci\'on lineal que recibe el nombre de ``funci\'on objetivo'', de manera que las variables de tal funci\'on est\'en sujetas a una serie de restricciones expresadas a trav\'es de un sistema de ecuaciones o inecuaciones lineales [26].
\end{itemize}

La investigaci\'on se centra en los m\'etodos heur\'isticos, los cuales se han caracterizado por ser procedimientos que garantizan buenas soluciones mediante una exploraci\'on limitada del espacio de b\'usqueda. Los m\'etodos heur\'isticos se pueden clasificar en m\'etodos constructivos y heur\'isticas de mejora [15], [26]. Algunas heur\'isticas utilizadas en la literatura para solucionar el MDVRP son:

\begin{itemize}
  \item Heur\'istica ad hoc: en 1993 fue presentada una nueva heur\'istica ad hoc en el aporte de [34], su m\'etodo de soluci\'on asigna los clientes a los dep\'ositos de manera f\'acil y r\'apida. En el procedimiento de mejora, permite que la distancia total recorrida por la flota aumente con la esperanza de encontrar finalmente rutas muy cortas [26].
  \item Algoritmo de Ahorros: extensi\'on del m\'etodo cl\'asico para VRP al contexto de m\'ultiples dep\'ositos. Asigna clientes a dep\'ositos y luego construye rutas fusionando recorridos seg\'un el ahorro de distancia [3]. En [35] se present\'o la primera adaptaci\'on para MDVRP.
  \item Algoritmo de Barrido: m\'etodo de dos fases: primero agrupa clientes por \'angulo polar respecto al dep\'osito, luego construye rutas para cada grupo. Bueno para problemas de tama\~no medio-grande [36].
  \item \textit{Merge-Head and Drop-Tail}: una t\'ecnica planteada para resolver el MDVRP es la denominada estrategia \textit{Merge-Head and Drop-Tail} (MHDT), combinada con optimizaci\'on de combinaci\'on de sectores (SCO), aplicada por [37]; aqu\'i se plantea un comportamiento inverso a la entrega, es decir debe generar rutas de recolecci\'on, dada esta \'unica diferencia, el resto de las caracter\'isticas y restricciones son id\'enticas, su propuesta funciona de igual manera para el MDVRP. En este m\'etodo, primero asignan puntos de recolecci\'on de acuerdo con la distancia m\'as cercana a los puntos de recolecci\'on, luego cada planta resuelve el VRP respectivamente, asignando clientes a los veh\'iculos y planificando el orden en que los veh\'iculos visitan a los clientes. En el \'ultimo paso, proponen el algoritmo de optimizaci\'on de combinaci\'on de sectores (SCO) para generar m\'ultiples soluciones iniciales, y luego estas soluciones iniciales se mejoran utilizando la estrategia de combinaci\'on de cabeza y ca\'ida (MHDT).
\end{itemize}

Las metaheur\'isticas son estrategias de alto nivel que gu\'ian y modifican heur\'isticas subyacentes para explorar el espacio de b\'usqueda de forma eficiente. Estos algoritmos se basan en la exploraci\'on iterativa del espacio de soluciones, utilizando estrategias para escapar de los \'optimos locales. Estas se clasifican en basadas en poblaci\'on o basadas en un punto [38].

La B\'usqueda Tab\'u se mantiene como uno de los mejores enfoques para resolver el MDVRP. Esta es una metaheur\'istica que explora el espacio de soluciones permitiendo movimientos que empeoran la soluci\'on actual para escapar de \'optimos locales. Para evitar ciclos, mantiene una lista de movimientos prohibidos o tab\'u durante un n\'umero determinado de iteraciones [33].

Tambi\'en otra de las metaheur\'isticas de punto \'unico que se destaca es la B\'usqueda Local Iterada. Este algoritmo parte de una soluci\'on inicial, aplica una perturbaci\'on para diversificar la b\'usqueda, luego una b\'usqueda local para alcanzar un nuevo \'optimo y finalmente decide si acepta la nueva soluci\'on seg\'un alg\'un criterio [26].

Entre los poblacionales se destacan bastante los gen\'eticos por la flexibilidad ante ciertas restricciones. Estos algoritmos mantienen una poblaci\'on de soluciones que evoluciona mediante operadores de selecci\'on, cruce y mutaci\'on. Para el MDVRP, las soluciones se representan como secuencias de clientes. El cruce combina sub-rutas de dos padres para generar un descendiente, reparando la factibilidad si es necesario [33].

Otro de los algoritmos poblacionales que se destacan fue el Algoritmo Colonia de Hormigas. Esta simula el comportamiento de hormigas reales que depositan feromona en el camino. Las hormigas artificiales construyen soluciones probabil\'isticamente guiadas por la feromona (que refleja la calidad de soluciones previas) y por informaci\'on heur\'istica local. Para MDVRP, se optimizan rutas por dep\'osito y se a\~naden operadores de cruce para mejorar el rendimiento [26].

\section{Algoritmos de agrupamiento para la fase de asignaci\'on de clientes a dep\'ositos}

La miner\'ia de datos es un campo interdisciplinario de la computaci\'on que se centra en el descubrimiento de patrones significativos en grandes vol\'umenes de datos mediante el uso de m\'etodos avanzados de inteligencia artificial, aprendizaje autom\'atico, estad\'istica y bases de datos. Su objetivo principal es extraer conocimiento valioso de los datos y transformarlo en una estructura comprensible, ya sea para su an\'alisis o para su almacenamiento en bases de datos con miras a usos posteriores [39].

La miner\'ia de datos se basa en una variedad de t\'ecnicas y m\'etodos que permiten extraer informaci\'on \'util de grandes vol\'umenes de datos. Algunas de las t\'ecnicas m\'as comunes incluyen [40]:

\begin{itemize}
  \item Clasificaci\'on: es una t\'ecnica utilizada para asignar elementos de un conjunto de datos a categor\'ias predefinidas. Esto se logra mediante algoritmos de aprendizaje supervisado, como los \'arboles de decisi\'on, las m\'aquinas de vectores de soporte (SVM) y las redes neuronales.
  \item Agrupamiento: es una t\'ecnica no supervisada que organiza los datos en cl\'usteres basados en la similitud entre los datos. Esta no requiere categor\'ias predefinidas.
  \item Regresi\'on: es una t\'ecnica utilizada para predecir un valor num\'erico basado en la relaci\'on entre variables.
  \item Reglas de asociaci\'on: identifican relaciones entre variables en un conjunto de datos.
\end{itemize}

La fase de asignaci\'on de clientes a dep\'ositos en MDVRP puede considerarse un problema de agrupamiento, ya que el resultado es un conjunto de clientes agrupados por dep\'osito. Esta interpretaci\'on sugiere un enfoque que agrupa a los clientes y luego programa los veh\'iculos en cada grupo. En la literatura el MDVRP puede resolverse en dos etapas: primero, los clientes deben asignarse a los dep\'ositos; luego, se deben construir rutas que vinculen a los clientes asignados al mismo dep\'osito. Un enfoque razonable ser\'ia dividir el problema en tantos subproblemas como dep\'ositos existan y resolver cada subproblema por separado, lo cual transforma el MDVRP en m\'ultiples VRP independientes [41]. Por tanto, resulta interesante adaptar los algoritmos de agrupamiento para la fase de asignaci\'on en el MDVRP.

Los algoritmos de agrupamiento se basan en la cuantificaci\'on de qu\'e tan similares o no sean los objetos del conjunto de datos. Las m\'etricas de distancia desempe\~nan un papel central en la agrupaci\'on de estos objetos en grupos. Entre las medidas comunes de similitud y distancia se incluyen varios tipos de distancia como \textit{Euclidian}, \textit{Manhattan}, \textit{Haversine} y \textit{Chebyshev}. La elecci\'on de la medida de proximidad influye en los resultados de la agrupaci\'on, ya que los distintos conjuntos de datos requieren t\'ecnicas de separaci\'on espec\'ificas y m\'etricas de similitud adaptadas a sus caracter\'isticas. Estos algoritmos de agrupamiento seg\'un [42], se pueden clasificar como:

\begin{itemize}
  \item Agrupamiento particional: en estos algoritmos cada partici\'on representa un grupo y cada objeto debe pertenecer exactamente a un grupo. Si se aplica la partici\'on difusa, permite que los objetos pertenezcan a m\'ultiples grupos con probabilidades. Aqu\'i se encuentran algoritmos como k-means, k-medoids, k-modes y c-means difusos. Estos m\'etodos son sensibles a la inicializaci\'on y a los valores at\'ipicos, y se obtienen resultados confiables cuando los grupos son esf\'ericos, de tama\~no y distribuci\'on uniformes. En la Figura~\ref{fig:particional} se puede apreciar gr\'aficamente el funcionamiento de este tipo de agrupamientos.
  \item Agrupamiento jer\'arquico: se dividen en enfoques basados en aglomeraci\'on (de abajo hacia arriba) y divisivos (de arriba hacia abajo), que producen una jerarqu\'ia de grupos representados como dendrogramas. Los m\'etodos basados en aglomeraci\'on comienzan con cada punto de datos como un grupo separado y fusionan los grupos m\'as cercanos de forma iterativa, mientras que los m\'etodos divisivos comienzan con todos los puntos de datos en un grupo y los dividen de forma recursiva. Los criterios de vinculaci\'on incluyen la vinculaci\'on simple, completa, promedio y el m\'etodo de Ward, cada uno de los cuales define la proximidad del grupo de manera diferente. Los ejemplos de algoritmos de agrupamiento jer\'arquico incluyen BIRCH, CURE y ROCK. En la Figura~\ref{fig:jerarquico} se puede apreciar gr\'aficamente el funcionamiento de este tipo de agrupamientos.
  \item Agrupamiento basado en densidad: los algoritmos de agrupamiento basados en densidad identifican los agrupamientos como regiones de alta densidad separadas por \'areas de baja densidad, sin requerir el n\'umero de agrupamientos como entrada. Dentro de esta clasificaci\'on se encuentra el agrupamiento espacial basado en densidad de aplicaciones con ruido (DBSCAN), el agrupamiento de puntos para identificar la estructura del agrupamiento (OPTICS) y el agrupamiento de picos de densidad (DPC). Estos m\'etodos pueden presentar dificultades con densidades variables. En la Figura~\ref{fig:densidad} se puede apreciar gr\'aficamente el funcionamiento de este tipo de agrupamientos.
\end{itemize}

\figura{fig03_particional.png}{0.5\textwidth}{Representaci\'on gr\'afica de algoritmos de agrupamiento particionales [43].}{fig:particional}

\figura{fig04_jerarquico.png}{0.5\textwidth}{Representaci\'on gr\'afica de algoritmos de agrupamiento jer\'arquicos [43].}{fig:jerarquico}

\figura{fig05_densidad.png}{0.5\textwidth}{Representaci\'on gr\'afica de algoritmos de agrupamiento basados en densidad [43].}{fig:densidad}

Los algoritmos de agrupamiento se aplican al MDVRP principalmente bajo el enfoque ``agrupar primero, rutear despu\'es'', donde primero se agrupan clientes y luego se construyen rutas para cada grupo. Un buen ejemplo de la aplicaci\'on de esto es una distribuci\'on de alimentos a supermercados en Reino Unido donde se utiliz\'o K-means para agrupar 120 supermercados en cuatro cl\'usteres correspondientes a los dep\'ositos disponibles, minimizando la distancia intra-cluster. Luego se optimizaron rutas independientes para cada cl\'uster. El ahorro en costos operativos fue del 15\% [43].

\section{Biblioteca de Heur\'isticas de Asignaci\'on en la Planificaci\'on de Rutas de Veh\'iculos con M\'ultiples Dep\'ositos}

El \'area de optimizaci\'on ha llevado a la creaci\'on de diversas bibliotecas en diferentes lenguajes de programaci\'on, cada una con su propia variedad de algoritmos y caracter\'isticas. A continuaci\'on, se presenta una lista no exhaustiva de bibliotecas destacadas en resolver VRP y sus variantes, entre las cuales se incluye el MDVRP. De cada una se ofrece una breve descripci\'on de sus caracter\'isticas y usos [15]:

\begin{itemize}
  \item JSprit: es una biblioteca de c\'odigo abierto escrita en Java. Es un componente ligero y flexible que resulta f\'acil de utilizar, fundamentado en una \'unica metaheur\'istica multiprop\'osito que actualmente aborda una variedad amplia de problemas, como el Problema del Agente Viajero y sus variantes de VRP. Aunque no est\'a espec\'ificamente dise\~nada para abordar la fase de asignaci\'on en problemas con m\'ultiples dep\'ositos, JSprit puede implementar estrategias de agrupamiento geogr\'afico como un paso preliminar, utilizando t\'ecnicas como ordenamientos angulares o agrupaciones espaciales similares al enfoque del cl\'asico algoritmo de barrido. Estas estrategias permiten organizar a los clientes en funci\'on de su proximidad antes de aplicar su metaheur\'istica interna [44].
  \item \textit{Google Operation Research Tools} (OR-Tools): es un software de c\'odigo abierto desarrollado en C++, con versiones estables disponibles en Java, C\# y Python. Originalmente dise\~nado para abordar diversos problemas de optimizaci\'on como la programaci\'on lineal, la programaci\'on entera y la programaci\'on con restricciones, esta herramienta ha ido incorporando gradualmente capacidades para resolver variantes del VRP. Aunque brinda herramientas para construir modelos completos de VRP, su arquitectura flexible permite personalizar fases espec\'ificas, como la asignaci\'on, mediante programaci\'on y heur\'isticas personalizadas. Adem\'as, se pueden integrar algoritmos de agrupamiento para realizar una asignaci\'on previa de clientes a los dep\'ositos. Tambi\'en es posible implementar restricciones a medida en el modelo, como l\'imites de capacidad o balanceo de carga, para ajustar la asignaci\'on inicial antes de continuar con la planificaci\'on de las rutas de los veh\'iculos. Esta herramienta resuelve varias variantes de VRP de manera nativa, pero no el MDVRP [45].
  \item \textit{Vehicle Routing Open-Source Optimization Machine} (VROOM): es un componente de software para problemas de planificaci\'on de rutas de veh\'iculos desarrollado en el lenguaje de programaci\'on C++. Esta herramienta utiliza \textit{Open Street Map} (OSM) y \textit{Open Source Routing Machine} (OSRM) como \textit{backend} para obtener rutas y devolver soluciones para las diferentes variantes del VRP. VROOM puede resolver distintas variantes del VRP que incluyen m\'ultiples dep\'ositos, capacidades, ventanas de tiempo, as\'i como entrega y recogida [45].
  \item Scikit-learn: es una biblioteca de aprendizaje autom\'atico en Python que, aunque no est\'a dise\~nada espec\'ificamente para resolver problemas de optimizaci\'on como el MDVRP, proporciona una amplia gama de algoritmos de agrupamiento y an\'alisis de datos que pueden ser utilizados para abordar la fase de asignaci\'on en este tipo de problemas. Algoritmos como K-means, DBSCAN y el agrupamiento jer\'arquico son \'utiles para preprocesar los datos y asignar clientes a dep\'ositos de manera eficiente, bas\'andose en proximidad y restricciones de capacidad. Su flexibilidad y compatibilidad con otras bibliotecas de Python la convierten en una herramienta vers\'atil para integrar t\'ecnicas de agrupamiento en soluciones personalizadas para el MDVRP [46].
\end{itemize}

En el \'ambito nacional, se han desarrollado diversas soluciones en el contexto acad\'emico orientadas a la resoluci\'on de problemas de planificaci\'on de rutas de veh\'iculos, como parte del proyecto de Optimizaci\'on y Metaheur\'isticas de la Facultad de Ingenier\'ia Inform\'atica de la Universidad Tecnol\'ogica de La Habana ``Jos\'e Antonio Echeverr\'ia'' (CUJAE). Algunas de estas herramientas se caracterizan por su enfoque modular y por estar dise\~nadas para abordar distintas variantes del VRP, con especial \'enfasis en la utilizaci\'on de heur\'isticas y metaheur\'isticas aplicadas al MDVRP [15]:

\begin{itemize}
  \item Bibliotecas de Heur\'isticas de Construcci\'on para Problemas de Planificaci\'on de Rutas de Veh\'iculos (BHCVRP): es un componente que implementa diez heur\'isticas de construcci\'on para resolver siete variantes del VRP, incluyendo el MDVRP. Actualmente las soluciones obtenidas mediante estas heur\'isticas pueden ser utilizadas como soluci\'on final o como parte de una soluci\'on m\'as compleja en un contexto donde se empleen metaheur\'isticas [47].
  \item Capa de servicio para soluciones a Problemas de Planificaci\'on de Rutas de Veh\'iculos (CaSVRP): es una capa de servicios web dirigidos a la resoluci\'on de diferentes variantes de VRP mediante algoritmos metaheur\'isticos. Entre las variantes del VRP que resuelve se halla MDVRP [48]. Esta soluci\'on se compone de un m\'odulo de optimizaci\'on donde se utiliza la Biblioteca de Clases para la Integraci\'on de Algoritmos Metaheur\'isticos (BiCIAM) [49], un m\'odulo para la generaci\'on de soluciones iniciales empleando BHCVRP [50] y otro para la generaci\'on de nuevas soluciones a partir de otras soluciones haciendo uso de LMOVRP [51].
  \item Biblioteca de Heur\'isticas de Asignaci\'on para Problemas de Planificaci\'on de Rutas de Veh\'iculos con M\'ultiples Dep\'ositos (BHAVRP): es una biblioteca que resuelve la fase de asignaci\'on de clientes a dep\'ositos en el MDVRP. Esta herramienta contiene 18 algoritmos de agrupamiento y heur\'isticas cl\'asicas de asignaci\'on, y es utilizada por otras herramientas como BHCVRP, CaSVRP y TransO [15].
\end{itemize}

La Tabla~\ref{tab:bibliotecas} presenta una comparaci\'on entre las bibliotecas antes mencionadas que resuelven la fase de asignaci\'on del MDVRP. Es importante resaltar criterios como el hecho de ser de c\'odigo abierto ya que eso ser\'ia de utilidad para seguir mejorando el componente.

\begin{table}[H]
\centering
\caption{Comparaci\'on entre las herramientas que resuelven el MDVRP.}
\label{tab:bibliotecas}
\small
\begin{tabular}{@{}lcccc@{}}
\toprule
\textbf{Biblioteca} & \textbf{\shortstack{Implementan\\heur\'isticas cl\'asicas}} & \textbf{\shortstack{Implementan\\algoritmos de\\agrupamiento}} & \textbf{\shortstack{Lenguajes\\soportados}} & \textbf{\shortstack{C\'odigo\\abierto}} \\
\midrule
JSprit & No & No & Java & S\'i \\
OR-Tools & S\'i & No & \shortstack{C++, Python,\\Java, C\#} & S\'i \\
VROOM & No & No & C++ & No \\
Scikit-learn & No & S\'i & Python & S\'i \\
CaSVRP & S\'i & No & Java & S\'i \\
BHCVRP & S\'i & No & \shortstack{Java, Python} & S\'i \\
BHAVRP & S\'i & S\'i & \shortstack{Java, Python} & S\'i \\
\bottomrule
\end{tabular}
\end{table}

A partir del an\'alisis comparativo de las bibliotecas podemos ver que solo BHAVRP y Scikit-learn ofrecen algoritmos de agrupamiento, lo que los posiciona como herramientas especialmente \'utiles cuando las necesidades incluyen esta capacidad para el preprocesamiento o modelado del problema. BHAVRP es la \'unica biblioteca espec\'ificamente enfocada en resolver la fase de asignaci\'on del MDVRP, mediante la implementaci\'on de heur\'isticas cl\'asicas y algoritmos de agrupamiento. Las dem\'as bibliotecas analizadas abordan el problema completo o incluyen la asignaci\'on como un paso preliminar, sin especializaci\'on expl\'icita en esta etapa.

BHAVRP actualmente ha sido integrado como parte funcional en herramientas como: CaSVRP y BHCVRP. Aunque su uso se ha concentrado en estos sistemas, su arquitectura modular y sus implementaciones en Java y Python permiten su potencial integraci\'on en otras aplicaciones compatibles con estos entornos [15].

Como se puede apreciar en la Figura~\ref{fig:flujo}, el flujo de ejecuci\'on comienza en el software que consume BHAVRP, donde se introducen los datos correspondientes a una instancia del problema MDVRP. Estos datos se env\'ian al componente BHAVRP, el cual procede a cargar los datos del problema y a determinar el tipo de distancia que se utilizar\'a en el c\'alculo de costos. A continuaci\'on, BHAVRP verifica si los datos fueron cargados correctamente. Si los datos presentan errores, el proceso se detiene. En caso contrario, se configuran los par\'ametros de ejecuci\'on necesarios, y se genera la matriz de costo correspondiente seg\'un la distancia seleccionada. Una vez generada la matriz de costos, el algoritmo se ejecuta en el contexto de BHAVRP, utilizando los par\'ametros definidos y la matriz de costos generada previamente. Al finalizar la ejecuci\'on del algoritmo, BHAVRP genera una soluci\'on, que es devuelta al sistema solicitante. Este \'ultimo obtiene as\'i los resultados de la asignaci\'on de clientes a dep\'ositos, completando el ciclo de interacci\'on [15].

\figura{fig06_flujo_bhavrp.png}{0.6\textwidth}{Diagrama de actividades del flujo de BHAVRP.}{fig:flujo}

Para soportar la variedad de algoritmos y garantizar su correcta ejecuci\'on, BHAVRP presenta una arquitectura n-capas con un enfoque basado en reutilizaci\'on, promoviendo consistencia con los lenguajes de desarrollo seleccionados [15]. Esta arquitectura se presenta en la Figura~\ref{fig:arquitectura}, que ofrece una vista detallada de BHAVRP.

\figura{fig07_arquitectura.png}{0.85\textwidth}{Diagrama de arquitectura n-capas con enfoque en reutilizaci\'on de BHAVRP.}{fig:arquitectura}

La capa espec\'ifica contiene los elementos responsables de gestionar la ejecuci\'on del sistema, como la clase Controller, que coordina la carga de datos, la configuraci\'on de par\'ametros y el flujo general de ejecuci\'on. Por su parte, en la capa general, que representa el n\'ucleo de la biblioteca es responsable de la definici\'on del problema, la modelaci\'on de las instancias y la ejecuci\'on de los algoritmos de asignaci\'on. Por su parte, la biblioteca trabaja con una estructura de datos que representan los elementos clave del MDVRP: clientes, dep\'ositos, y sus respectivas restricciones y par\'ametros. La estructura b\'asica para utilizar la biblioteca incluye un encabezado con el n\'umero de dep\'ositos, clientes y veh\'iculos disponibles por dep\'osito, seguido de las capacidades m\'aximas de los veh\'iculos. Los clientes y los dep\'ositos est\'an definidos por un identificador y sus coordenadas (x, y). Adem\'as, cada cliente tiene una demanda asociada, mientras que cada dep\'osito cuenta con una capacidad equivalente a la suma de las capacidades de los veh\'iculos en su flota. BHAVRP en la capa intermedia contiene varias herramientas. Una de estas permite trabajar tanto con distancias reales como aproximadas (Euclidean, Manhattan, Chebyshev y Haversine). Luego para abordar el c\'alculo aproximado de las distancias, integra la herramienta OSRM (Open Source Routing Machine), que permite calcular distancias y tiempos de viaje basados en mapas reales, tanto en modalidad local como a trav\'es de solicitudes a un servidor en l\'inea. En la capa de software se implementa la comunicaci\'on de la herramienta OSRM que expone sus servicios a trav\'es de una API HTTP RESTful. El cliente Java (como osrm-java-client) construye las peticiones HTTP y convierte las respuestas a JSON [15].

\section{Conclusiones parciales}

A partir del estudio realizado sobre el marco te\'orico de los problemas de planificaci\'on de rutas de veh\'iculos, espec\'ificamente la variante con m\'ultiples dep\'ositos, as\'i como los m\'etodos de soluci\'on y algoritmos de Miner\'ia de Datos existentes para resolver estos problemas se puede concluir que:

\begin{itemize}
  \item Los VRP son problemas de optimizaci\'on combinatoria clasificados como NP-Duro por lo que se utilizan principalmente algoritmos heur\'isticos y metaheur\'isticos para su soluci\'on. El MDVRP es una de las numerosas variantes del VRP y requiere otro tipo de estrategias de soluci\'on.
  \item Una de las estrategias para la soluci\'on del MDVRP es dividir en la fase de asignaci\'on y la fase de planificaci\'on de las rutas. Adem\'as, para esta primera fase de asignaci\'on resulta efectivo el uso de algoritmos de miner\'ia de datos.
  \item BHAVRP es una biblioteca que permite resolver esa primera fase de asignaci\'on y tiene la facilidad para a\~nadir nuevas heur\'isticas y algoritmos de agrupamiento gracias a su arquitectura n-capas orientada a reutilizaci\'on.
  \item La incorporaci\'on de un algoritmo basado en densidad resulta interesante ya que presenta un enfoque diferente a los algoritmos particionales y jer\'arquicos que ya presenta BHAVRP.
  \item La incorporaci\'on de otro algoritmo de agrupamiento jer\'arquico a la biblioteca enriquece la biblioteca con otro algoritmo de esta categor\'ia y verifica si UPGMC presenta los mejores resultados gracias a su comportamiento jer\'arquico o a su naturaleza intr\'inseca.
\end{itemize}
%==============================================================
%                    CAPITULO 2
%==============================================================
\chapter{Incorporaci\'on de algoritmos a la Biblioteca de Heur\'isticas de Asignaci\'on del Problema de Planificaci\'on de Rutas con M\'ultiples Dep\'ositos}

En este cap\'itulo se presenta la incorporaci\'on de nuevos algoritmos a BHAVRP. Se muestra una tabla comparativa entre los algoritmos de agrupamiento ya implementados en la biblioteca y los que se proponen. Adem\'as, se presenta una descripci\'on detallada de los algoritmos de agrupamiento propuestos, modificaciones en la arquitectura de la biblioteca para la incorporaci\'on de estos y la adaptaci\'on para el MDVRP mediante diagramas de actividades.

\section{Algoritmo DBSCAN}

El algoritmo \textit{Density-Based Spatial clustering of Applications with Noise}, pertenece a la familia de algoritmos de agrupamiento basado en densidad, empleado frecuentemente cuando los datos contienen ruido. Este tipo de algoritmos se basan en el concepto de densidad de la vecindad de un punto, donde se calcula un n\'umero de puntos que est\'en dentro de un radio llamado \'Epsilon [52].

Para el funcionamiento del algoritmo se le pasan los siguientes par\'ametros [53]:

\begin{itemize}
  \item \'Epsilon ($\varepsilon$): el radio de la vecindad alrededor de un punto de datos.
  \item Puntos m\'inimos (\textit{minPts}): el n\'umero m\'inimo de puntos de datos que deben estar dentro del radio $\varepsilon$ de un punto dado.
\end{itemize}

Adem\'as, DBSCAN clasifica los puntos de los datos de la forma siguiente [52]:

\begin{enumerate}
  \item Punto n\'ucleo: un punto $P$ es un punto n\'ucleo si al menos \textit{minPts} est\'an a una distancia $\varepsilon$ de \'el o, esos puntos son directamente alcanzables desde $P$. No es posible tener puntos directamente alcanzables desde un punto que no sea un n\'ucleo.
  \item Frontera: un punto $Q$ es frontera si est\'a dentro del radio $\varepsilon$ de un punto n\'ucleo $P$, pero no tiene suficientes puntos vecinos dentro de su propio radio $\varepsilon$ para ser considerado n\'ucleo. Estos puntos pertenecen al cl\'uster asociado a $P$, pero no pueden expandir el grupo hacia nuevos puntos.
  \item Ruido: un punto que no sea alcanzable desde cualquier otro punto.
\end{enumerate}

Si $P$ es un punto n\'ucleo, entonces forma un grupo junto a otros puntos (n\'ucleo o no) que sean alcanzables desde \'el. Cada grupo contiene al menos un punto n\'ucleo. Los puntos no n\'ucleo alcanzables pueden pertenecer al grupo, pero act\'uan como frontera puesto que no es posible alcanzar m\'as puntos desde estos. Por lo tanto, el tiempo de respuesta se ve comprometido cuando se aplica en el contexto de grandes vol\'umenes de datos [52].

En la Figura~\ref{fig:dbscan_puntos} se puede apreciar de forma gr\'afica c\'omo funciona el algoritmo.

\figura{fig08_dbscan.png}{0.6\textwidth}{Ejemplo de la agrupaci\'on de puntos por densidad de DBSCAN.}{fig:dbscan_puntos}

Para encontrar un cl\'uster, DBSCAN parte de un punto arbitrario cualquiera y recupera todos los puntos alcanzables por su densidad con respecto a su radio $\varepsilon$ y los puntos m\'nimos. Si ese punto es un n\'ucleo, este procedimiento genera un cl\'uster. Si es un punto frontera, DBSCAN se sit\'ua en el siguiente punto de la base de datos. Dado que se utilizan valores globales para $\varepsilon$ y \textit{minPts}, DBSCAN puede fusionar dos cl\'usteres en un solo cl\'uster si dos cl\'usteres de diferente densidad est\'an cerca entre s\'i. Entonces, dos conjuntos de puntos con al menos la densidad del cl\'uster m\'as delgado solo se separar\'an si la distancia entre los dos conjuntos es mayor que el radio $\varepsilon$ [54].

Si dos cl\'usteres $C_{1}$ y $C_{2}$ est\'an muy cerca, podr\'ia ocurrir que alg\'un punto $P$ pertenezca a ambos. Entonces, $P$ debe ser un punto l\'imite en ambos cl\'usteres; de lo contrario, $C_{1}$ ser\'ia igual a $C_{2}$, dado que se utilizan par\'ametros globales. En este caso, el punto $P$ se asignar\'a al cl\'uster descubierto primero. Excepto en estas raras situaciones, el resultado de DBSCAN no es independiente del orden en que se visitan los puntos de la base de datos [54].

Con el objetivo de resumir el flujo descrito anteriormente, la Figura~\ref{fig:dbscan_flujo} presenta el diagrama de flujo del algoritmo DBSCAN en su versi\'on m\'as b\'asica. Este diagrama resume las etapas fundamentales del proceso: selecci\'on de puntos no visitados, verificaci\'on de densidad mediante los par\'ametros $\varepsilon$ y \textit{minPts}, expansi\'on de cl\'usteres y clasificaci\'on de puntos como ruido cuando no cumplen el criterio de densidad.

\figura{fig09_flujo_dbscan.png}{0.55\textwidth}{Diagrama de flujo del algoritmo DBSCAN.}{fig:dbscan_flujo}

DBSCAN fue concebido como un algoritmo de agrupamiento basado en densidad espacial, por lo que necesita adaptaciones para resolver MDVRP debido a la presencia de restricciones adicionales como la capacidad de los dep\'ositos y la necesidad de asignaci\'on directa hacia dep\'ositos. Por esta raz\'on, se desarrolla una adaptaci\'on del algoritmo que incorpora restricciones propias del MDVRP sin perder las funciones originales del enfoque basado en densidad.

La Figura~\ref{fig:dbscan_mdvrp} muestra el diagrama de actividades correspondiente a la versi\'on adaptada del algoritmo DBSCAN para el MDVRP.

\figura{fig10_flujo_dbscan_mdvrp.png}{0.6\textwidth}{Diagrama de flujo del algoritmo DBSCAN a MDVRP.}{fig:dbscan_mdvrp}

En esta variante, los cl\'usteres no se crean libremente desde cualquier punto n\'ucleo, sino que se inicializan expl\'icitamente a partir de los dep\'ositos existentes. A partir de cada dep\'osito, el algoritmo expande el cl\'uster considerando los clientes dentro del radio $\varepsilon$, siempre que se cumpla simult\'aneamente el criterio de densidad (\textit{minPts}) y la restricci\'on de capacidad del dep\'osito.

Adem\'as, se incorpora una fase final de reasignaci\'on, donde los clientes que no pudieron ser incluidos durante la expansi\'on inicial son evaluados nuevamente y asignados al dep\'osito m\'as cercano con capacidad disponible.

\section{Algoritmo CURE}

En [12] se presenta el algoritmo de agrupamiento jer\'arquico \textit{Clustering Using Representatives} que usa puntos representativos y un muestreo para gestionar conjuntos de datos m\'as grandes con menores costes computacionales. CURE se puede dividir en dos fases: inicializaci\'on y finalizaci\'on. La inicializaci\'on comienza tomando una muestra aleatoria sin reemplazo de la poblaci\'on original y aplicando el agrupamiento jer\'arquico. Tanto el n\'umero de puntos representativos como el factor de contracci\'on se pueden optimizar como parte del ajuste de par\'ametros. En la etapa de finalizaci\'on, para cada cl\'uster resultante, selecciona un peque\~no conjunto de puntos representativos bien dispersos dentro del cl\'uster y, tras la contracci\'on hacia el centroide del cl\'uster, se asigna el cl\'uster a todos los puntos restantes calculando las distancias hacia todos esos representantes del cl\'uster. Gracias a esto, el algoritmo puede incluso reconocer geometr\'ias de formas arbitrarias, a la vez que es bueno frente a valores at\'ipicos [55].

En la Figura~\ref{fig:cure_flujo_pasos} se muestra la secuencia de pasos que sigue el algoritmo [12]:

\figura{fig11_cure_pasos.png}{0.8\textwidth}{Flujo de CURE.}{fig:cure_flujo_pasos}

Para comprender el funcionamiento del algoritmo CURE, se mostrar\'a su proceso paso a paso, comenzando con el muestreo del conjunto de datos [56]:

\begin{itemize}
  \item Muestreo del conjunto de datos: CURE comienza seleccionando una muestra aleatoria representativa del conjunto de datos. El proceso de muestreo reduce el n\'umero de puntos de datos, lo que acelera la velocidad de c\'alculo y mantiene la integridad del cl\'uster. El muestreo aleatorio simple permite una implementaci\'on r\'apida, pero no logra capturar casos extremos cruciales. Las estrategias de muestreo implementadas por CURE permiten la agrupaci\'on posterior para representar con precisi\'on la complejidad y diversidad de todo el conjunto de datos.
  \item Partici\'on de cl\'usteres: tras muestrear el conjunto de datos, CURE aplica un m\'etodo jer\'arquico que divide los datos en subconjuntos manejables. El algoritmo mide las similitudes de los puntos de datos mediante m\'etricas de distancia, algunas de estas son la distancia euclidiana o la distancia de Manhattan. Las m\'etricas de c\'alculo de distancia son esenciales para determinar la proximidad de los puntos y establecer una base s\'olida para una fusi\'on de cl\'usteres eficaz. Al inicio del proceso, cada punto de datos funciona como su propio cl\'uster para representar la granularidad fina de los datos. El algoritmo realiza fusiones de cl\'usteres sucesivas utilizando est\'andares de proximidad para agrupar puntos similares. El algoritmo contin\'ua con las operaciones de fusi\'on de cl\'usteres hasta alcanzar un n\'umero definido de cl\'usteres o hasta que se activa una condici\'on de terminaci\'on alternativa. El proceso de selecci\'on de grupos en CURE crea cl\'usteres que se alinean con las categor\'ias naturales presentes en los datos. Mediante la partici\'on jer\'arquica, CURE supera las limitaciones de los algoritmos de agrupamiento tradicionales, que requieren que los cl\'usteres tengan formas convexas.
  \item Selecci\'on de puntos representativos: CURE selecciona m\'ultiples puntos representativos para representar cada cl\'uster, en lugar de depender de un \'unico centroide. Estos puntos seleccionados del cl\'uster capturan su rango espacial y estructura. CURE logra un mejor reconocimiento de los l\'imites del cl\'uster y una mejor comprensi\'on de la estructura interna al utilizar m\'ultiples puntos para representar cada cl\'uster. Despu\'es de seleccionar los puntos representativos, el algoritmo los mueve hacia la media del cl\'uster con una reducci\'on espec\'ifica. Esta reducci\'on hace que el algoritmo sea menos reactivo a los valores at\'ipicos, ya que mueve los puntos distantes hacia los puntos centrales del cl\'uster. El \'exito del algoritmo depende en gran medida de la cantidad de puntos representativos utilizados durante la ejecuci\'on. Usar pocos puntos representativos puede no capturar la complejidad del cl\'uster, mientras que usar demasiados puede aumentar los costos computacionales.
  \item Fusi\'on de cl\'usteres: tras identificar los puntos representativos, el algoritmo fusiona los cl\'usteres mediante un enfoque iterativo sistem\'atico. El procedimiento se basa en la medici\'on de la distancia entre los puntos representativos de diferentes cl\'usteres. Se utilizan m\'etricas predefinidas, como la distancia euclidiana, para medir las distancias, eliminando as\'i conflictos a la hora de encontrar los cl\'usteres m\'as cercanos. El algoritmo identifica qu\'e pares de cl\'usteres tienen la menor distancia respecto a la separaci\'on de los puntos representativos durante cada paso de la evaluaci\'on. El algoritmo toma decisiones precisas sobre la fusi\'on de cl\'usteres, preservando las relaciones espaciales y los patrones naturales de alineaci\'on de datos dentro de los cl\'usteres. El proceso itera hasta que se alcanzan algunos cl\'usteres predeterminados o se alcanza otro criterio de terminaci\'on. El criterio de terminaci\'on depende de factores como la distancia m\'inima entre cl\'usteres o la similitud m\'axima permitida dentro de ellos. Esto garantiza que los cl\'usteres generados correspondan a las agrupaciones naturales de los datos y sean lo suficientemente flexibles como para capturar formas irregulares y complejas.
  \item Manejo de valores at\'ipicos: cuando existen valores at\'ipicos, los resultados se distorsionan porque causan formas incorrectas de los cl\'usteres y malinterpretan la estructura de los datos. El algoritmo CURE resuelve estas limitaciones utilizando m\'ultiples puntos representativos, que identifican con precisi\'on la forma y distribuci\'on reales de los cl\'usteres.
\end{itemize}

La complejidad computacional en el peor caso de este algoritmo de agrupamiento se puede demostrar como $O(n^{2}\log(n))$. [12] demuestra que cuando la dimensionalidad de los puntos de datos es peque\~na, la complejidad temporal se reduce a\'un m\'as a $O(n^{2})$. Dado que se requiere un espacio lineal, se deduce que la complejidad espacial de este algoritmo es $O(n)$.

Para sintetizar todo el funcionamiento descrito anteriormente, la Figura~\ref{fig:cure_flujo} presenta el diagrama de actividades del algoritmo CURE en su versi\'on base. Este diagrama resume el proceso jer\'arquico de fusi\'on de cl\'usteres, comenzando con cada punto como un cl\'uster independiente y continuando con la uni\'on iterativa de los cl\'usteres m\'as cercanos hasta alcanzar el n\'umero deseado $k$. Destac\'andose la selecci\'on de m\'ultiples puntos representativos por cl\'uster y su posterior contracci\'on hacia el centroide, elemento diferenciador de CURE frente a otros algoritmos jer\'arquicos tradicionales.

\figura{fig12_flujo_cure.png}{0.55\textwidth}{Diagrama de flujo del algoritmo CURE.}{fig:cure_flujo}

La aplicaci\'on directa de CURE al MDVRP no contempla restricciones fundamentales como la capacidad de los dep\'ositos ni la necesidad de que los cl\'usteres est\'en asociados expl\'icitamente a estos. Por ello, se desarroll\'o una adaptaci\'on del algoritmo que preserva la l\'ogica de representantes m\'ultiples y fusi\'on basada en proximidad, pero incorpora las restricciones pertinentes. La Figura~\ref{fig:cure_mdvrp} muestra el diagrama de actividades correspondiente a la versi\'on adaptada de CURE para el MDVRP.

\figura{fig13_flujo_cure_mdvrp.png}{0.6\textwidth}{Diagrama de flujo del algoritmo CURE adaptado al MDVRP.}{fig:cure_mdvrp}

En esta variante, los cl\'usteres iniciales se dividen en dos conjuntos diferenciados: cl\'usteres de clientes individuales y cl\'usteres asociados a dep\'ositos.

El proceso iterativo ya no busca reducir el n\'umero total de cl\'usteres hasta alcanzar un valor $k$, sino que se orienta a fusionar progresivamente clientes con cl\'usteres de dep\'ositos, siempre que:

\begin{itemize}
  \item La distancia m\'inima entre puntos representativos sea la menor disponible.
  \item La capacidad del dep\'osito permita absorber la demanda del cliente o subcl\'uster.
\end{itemize}

Tras cada fusi\'on, se actualizan los puntos representativos del cl\'uster resultante, manteniendo la parte m\'as importante y que caracteriza al algoritmo CURE original.

\section{Comparaci\'on de los algoritmos de agrupamiento}

BHAVRP implementa diferentes tipos de algoritmos adaptados para resolver espec\'ificamente la fase de asignaci\'on del MDVRP desde diferentes enfoques. Estos seg\'un la arquitectura de la biblioteca se dividen en heur\'isticos, de agrupamiento y una categor\'ia otros para los que no se corresponde con alguna de las anteriores. Esta investigaci\'on se centra en los algoritmos de agrupamiento. A continuaci\'on, se describen los presentes en la biblioteca [15]:

\begin{itemize}
  \item \textit{K-Means}: asigna seleccionando los clientes con la menor distancia a un dep\'osito.
  \item \textit{Partitional Around Medoids} (PAM): realiza diferentes asignaciones de los clientes a los dep\'ositos con el objetivo de encontrar la asignaci\'on que minimice el costo.
  \item \textit{Clustering Large Application} (CLARA): es una extensi\'on de PAM para grandes conjuntos de datos, seleccionando muestras representativas.
  \item \textit{Farthest-First}: asigna seleccionando como centroides los elementos con la mayor distancia entre ellos.
  \item \textit{Unweighted Pair Group Method with Centroid} (UPGMC): agrupa los cl\'usteres m\'as cercanos hasta que la cantidad de cl\'usteres sea equivalente al n\'umero de dep\'ositos.
\end{itemize}

En la versi\'on actual de BHAVRP, los algoritmos de agrupamiento implementados presentan limitaciones importantes que justifican la incorporaci\'on de nuevos algoritmos. En la Tabla~\ref{tab:comparacion_algoritmos} se muestra una comparaci\'on entre todos los algoritmos de agrupamiento de la biblioteca y los algoritmos que se proponen para esta.

\begin{table}[H]
\centering
\caption{Comparaci\'on de algoritmos de BHAVRP con los propuestos.}
\label{tab:comparacion_algoritmos}
\small
\begin{tabular}{@{}lccccccc@{}}
\toprule
\textbf{Caracter\'istica} & \textbf{KMEANS} & \textbf{PAM} & \textbf{CLARA} & \textbf{\shortstack{Farthest\\First}} & \textbf{UPGMC} & \textbf{DBSCAN} & \textbf{CURE} \\
\midrule
\shortstack[l]{Requiere especificar\\n\'umero de cl\'usteres} & X & X & X & X & X &  & X \\
\shortstack[l]{Detecta formas arbitrarias\\en la forma de los cl\'usteres} &  &  &  &  &  & X & X \\
\shortstack[l]{Escalable a miles\\de puntos} & X &  & X & X &  & X & X \\
\shortstack[l]{Buena identificaci\'on\\de valores at\'ipicos} &  & X & X &  &  & X & X \\
\bottomrule
\end{tabular}
\end{table}

Como se puede apreciar, los algoritmos implementados cuentan con ciertas limitaciones, principalmente ante la forma de los cl\'usteres. Los algoritmos DBSCAN y CURE cubren estas caracter\'isticas ausentes en los m\'etodos actuales. Adem\'as, a diferencia de UPGMC estos pueden detectar formas arbitrarias en vez de formas circulares e identificar bien los valores at\'ipicos lo cual mejorar\'ia la calidad de asignaci\'on con respecto a este. Implementar una adaptaci\'on de estos algoritmos para la fase de asignaci\'on de clientes a dep\'ositos en MDVRP lograr\'a identificar formas m\'as complejas de cl\'usteres que se adaptan mejor a la realidad geogr\'afica.

\section{Incorporaci\'on de los algoritmos propuestos a BHAVRP}

La incorporaci\'on de los nuevos algoritmos de agrupamiento a la biblioteca BHAVRP requiri\'o la definici\'on de las siguientes clases: AbstractDensity, DBSCAN, CURE y CureCluster. Estas clases se integran a la arquitectura existente respetando el patr\'on de dise\~no que organiza los algoritmos de agrupamiento dentro de la biblioteca, basado en el patr\'on \textit{Template Method}. En la Tabla~\ref{tab:clases} se puede apreciar las clases con una peque\~na descripci\'on de estas:

\begin{table}[H]
\centering
\caption{Clases implementadas para la incorporaci\'on de los nuevos algoritmos de agrupamiento.}
\label{tab:clases}
\small
\begin{tabular}{@{}lp{10cm}@{}}
\toprule
\textbf{Clase} & \textbf{Descripci\'on} \\
\midrule
\textbf{AbstractDensity} & Clase abstracta que concentra operaciones comunes para algoritmos basados en densidad, como el c\'alculo de distancias, la b\'usqueda de vecinos dentro de un radio $\varepsilon$ y la asignaci\'on de clientes a cl\'usteres. \\
\textbf{DBSCAN} & Implementaci\'on del algoritmo DBSCAN adaptado al MDVRP. Realiza la expansi\'on de cl\'usteres desde los dep\'ositos considerando densidad de clientes, l\'imites de capacidad y posterior reasignaci\'on de clientes no asignados. \\
\textbf{CURE} & Implementaci\'on del algoritmo jer\'arquico CURE adaptado al MDVRP. Ejecuta la asignaci\'on de clientes a dep\'ositos mediante la fusi\'on progresiva de cl\'usteres basada en la distancia entre puntos representativos. \\
\textbf{CureCluster} & Estructura utilizada por CURE para representar un cl\'uster durante el proceso de agrupamiento. Mantiene los clientes del cl\'uster, su demanda acumulada, el centroide y los puntos representativos utilizados para calcular distancias entre cl\'usteres. \\
\bottomrule
\end{tabular}
\end{table}

Para esta actualizaci\'on de la biblioteca fue necesario modificar el paquete \textit{Clustering} de la capa general, ya que todas estas clases fueron implementadas en este. Como se puede apreciar en la Figura~\ref{fig:arquitectura_actualizada}, en azul se muestra el paquete modificado:

\figura{fig14_arquitectura_actualizada.png}{0.85\textwidth}{Diagrama de arquitectura actualizado.}{fig:arquitectura_actualizada}

En BHAVRP, todos los algoritmos de asignaci\'on de clientes a dep\'ositos siguen una estructura com\'un que consta de tres fases: inicializaci\'on, asignaci\'on iterativa y finalizaci\'on. Independientemente de la heur\'istica espec\'ifica, el flujo general es siempre el mismo: se preparan las estructuras de datos, se asignan los clientes a los dep\'ositos respetando las restricciones del problema, y se construye la soluci\'on final. Sin embargo, los detalles de c\'omo se realiza la inicializaci\'on, c\'omo se asignan los clientes en cada iteraci\'on y c\'omo se determina la convergencia var\'ian de un algoritmo a otro.

En este contexto, el patr\'on \textit{Template Method} se utiliza para definir el esqueleto del algoritmo en un m\'etodo plantilla que establece la secuencia invariante de ejecuci\'on. En el dise\~no de BHAVRP, la clase abstracta AssignmentTemplate (que hereda de la clase base Heuristic) encapsula el flujo del algoritmo en su m\'etodo plantilla \textit{toClustering} (declarado en la interfaz IAssignment). Este m\'etodo define la secuencia de pasos invariantes:

\begin{enumerate}
  \item Inicializaci\'on (\textbf{initialize}): crea las estructuras necesarias (cl\'usteres, listas de clientes, centroides, etc.) y prepara el estado inicial del algoritmo.
  \item Asignaci\'on (\textbf{assign}): ejecuta el ciclo iterativo donde los clientes se asignan a los cl\'usteres utilizando la matriz de costos, se actualizan los representantes de los cl\'usteres y se verifica si se requiere una nueva iteraci\'on.
  \item Finalizaci\'on (\textbf{finish}): ensambla la soluci\'on final, clasificando los clientes no asignados y los cl\'usteres resultantes.
\end{enumerate}

DBSCAN y CURE respetan este patr\'on al implementar las operaciones abstractas definidas en la plantilla sin modificar la estructura general del flujo. En ambos casos, el m\'etodo \textit{to\_clustering} mantiene la secuencia establecida por la clase base, mientras que cada algoritmo redefine los m\'etodos \textit{initialize}, \textit{assign} y \textit{finish} para incorporar su l\'ogica particular de agrupamiento. De esta manera, DBSCAN implementa la expansi\'on de cl\'usteres basada en densidad y CURE el proceso jer\'arquico de fusi\'on mediante puntos representativos, pero ambos siguen el mismo esquema de ejecuci\'on definido por la biblioteca.

La incorporaci\'on de estas clases respeta varios principios fundamentales de dise\~no de software:

\begin{itemize}
  \item Principio Abierto/Cerrado: la arquitectura base de BHAVRP no requiere modificaciones para incorporar nuevos algoritmos de agrupamiento. La extensi\'on del sistema se logra mediante la creaci\'on de nuevas clases que implementan los m\'etodos definidos en la jerarqu\'ia existente.
  \item Principio de responsabilidad \'unica: cada clase mantiene una responsabilidad bien definida dentro del dise\~no. La estructura general del algoritmo se encuentra encapsulada en la clase base, mientras que las clases concretas se encargan exclusivamente de implementar la l\'ogica particular de cada t\'ecnica de agrupamiento.
  \item Principio de favorecer composici\'on sobre la herencia: en la clase CureCluster en lugar de extender directamente la clase Cluster mediante herencia, CureCluster encapsula una instancia de esta clase y a\~nade informaci\'on adicional necesaria para el funcionamiento del algoritmo CURE, como los puntos representativos, el centroide y la colecci\'on de clientes asociados al cl\'uster.
  \item Alta cohesi\'on: las operaciones relacionadas con el c\'alculo de distancias, vecindades y otras funcionalidades compartidas se agrupan en clases espec\'ificas como AbstractDensity. Esto permite que cada componente del sistema concentre funcionalidades relacionadas entre s\'i.
  \item Bajo acoplamiento: las implementaciones concretas dependen \'unicamente de las interfaces y estructuras generales proporcionadas por BHAVRP, lo que facilita la sustituci\'on o incorporaci\'on de nuevas heur\'isticas sin afectar el resto del sistema.
\end{itemize}

En la Figura~\ref{fig:clases} se puede apreciar la jerarqu\'ia de clases donde se aplica el patr\'on Template Method y la implementaci\'on de las nuevas clases respetando este. La clase AssignmentTemplate act\'ua como clase base con el m\'etodo plantilla, y las distintas heur\'isticas de asignaci\'on sobrescriben los m\'etodos abstractos seg\'un su l\'ogica espec\'ifica.

\figura{fig15_diagrama_clases.png}{0.9\textwidth}{Diagrama de clases actualizado.}{fig:clases}

\section{Conclusiones parciales}

A partir del an\'alisis sobre los algoritmos DBSCAN y CURE, la comparaci\'on con los algoritmos de agrupamiento de BHAVRP, adem\'as del dise\~no de la adaptaci\'on de estos algoritmos a MDVRP se concluye que:

\begin{itemize}
  \item Los algoritmos DBSCAN y CURE suponen una buena inclusi\'on a BHAVRP gracias a su capacidad de generar cl\'usteres con formas m\'as complejas y la detecci\'on de valores at\'ipicos. DBSCAN siendo un buen primer acercamiento a los algoritmos basados en densidad dentro de la biblioteca y CURE una alternativa bastante robusta gracias al uso de puntos representativos.
  \item DBSCAN, originalmente dise\~nado para detecci\'on de densidad, puede reinterpretarse como un mecanismo de expansi\'on controlada desde dep\'ositos, convirti\'endolos en n\'ucleos de agrupamiento.
  \item CURE, como algoritmo jer\'arquico, ofrece mayor flexibilidad al permitir fusiones progresivas basadas en distancia m\'inima entre puntos representativos.
  \item Ambas adaptaciones mantienen la esencia estructural de los algoritmos cl\'asicos, pero introducen mecanismos adicionales de validaci\'on que garantizan factibilidad log\'istica.
\end{itemize}
%==============================================================
%                    CAPITULO 3
%==============================================================
\chapter{An\'alisis experimental}

En este cap\'itulo se presentan las pruebas funcionales y unitarias para las funcionalidades principales de los algoritmos implementados, adem\'as de un estudio comparativo entre los resultados obtenidos por los algoritmos de agrupamiento de la biblioteca y los nuevos propuestos usando las m\'etricas de evaluaci\'on que contiene BHAVRP. Tambi\'en una prueba estad\'istica no param\'etrica entre los dos algoritmos jer\'arquicos.

Todos los experimentos se llevan a cabo en un entorno controlado, utilizando un ordenador con las siguientes especificaciones: procesador AMD A6-7310 APU with AMD Radeon R4 Graphics de 2.4GHz, 4 GB de memoria RAM y sistema operativo Linux Mint de 64 bits, usando el IDE Intellij IDEA con Java en su versi\'on 23.

\section{Pruebas funcionales}

Con el prop\'osito de garantizar la correcta integraci\'on y funcionamiento de los nuevos algoritmos implementados en BHAVRP, se desarrolla un conjunto de pruebas unitarias usando el componente Junit para validar las funciones esenciales de los algoritmos: m\'etodo \textit{distance} de la clase DBSCAN y constructor, as\'i como m\'etodo \textit{merge} de la clase CureCluster y el m\'etodo \textit{calculateClusterDistance} de la clase CURE. Estas pruebas funcionales se realizan previamente al an\'alisis de calidad y rendimiento de los algoritmos de asignaci\'on sobre instancias.

\subsection{Prueba del m\'etodo \textit{distance} en DBSCAN}

El m\'etodo \textit{distance} de la clase DBSCAN es fundamental para el algoritmo de agrupamiento basado en densidad, ya que calcula la distancia euclidiana entre dos puntos geogr\'aficos representados por objetos de la clase \textit{Location}. Para asegurar su correcto funcionamiento se dise\~nan pruebas unitarias que verifican:

\begin{itemize}
  \item C\'alculo correcto de distancias: se comparan distancias conocidas, como la distancia entre (0,0) y (3,4) que debe ser 5.0 por la m\'etrica euclidiana.
  \item Distancia a s\'i mismo: se verifica que la distancia de un punto a s\'i mismo sea cero, lo cual es una propiedad b\'asica de cualquier m\'etrica.
  \item Manejo de coordenadas negativas: se prueba con puntos que incluyen coordenadas negativas para asegurar que el c\'alculo no se ve afectado por el signo.
  \item Simetr\'ia: se comprueba que la distancia entre A y B sea igual a la distancia entre B y A, propiedad necesaria para cualquier funci\'on de distancia.
\end{itemize}

Adicionalmente, se verifica que el m\'etodo sea consistente para diferentes instancias y que el tipo de distancia por defecto sea la euclidiana. En la Figura~\ref{fig:prueba_distance} se puede apreciar el resultado proporcionado por el IDE a esta prueba.

\figura{fig16_prueba_distance.png}{0.85\textwidth}{Resultados de las pruebas del m\'etodo distance.}{fig:prueba_distance}

\subsection{Pruebas del constructor y m\'etodo merge de la clase CureCluster}

La clase CureCluster representa un cl\'uster en el algoritmo CURE. Su correcta construcci\'on y la operaci\'on de fusi\'on (\textit{merge}) son cr\'iticas para el proceso jer\'arquico. Las pruebas unitarias se centran en:

Creaci\'on de cl\'usteres:

\begin{itemize}
  \item Con un solo cliente: se verifica que el cl\'uster se cree correctamente, que contenga al cliente, que el centroide se calcule adecuadamente y que la demanda total sea la esperada.
  \item Con m\'ultiples clientes: se comprueba que el centroide sea el promedio de las coordenadas de los clientes y que la demanda sea la suma de las demandas individuales.
  \item Se verifica que los puntos representativos se generen correctamente tras la creaci\'on.
\end{itemize}

Operaci\'on de fusi\'on (merge):

\begin{itemize}
  \item Se fusionan dos cl\'usteres y se comprueba que el n\'umero de clientes resultante sea la suma de ambos, que la demanda total sea la suma y que el centroide se actualice correctamente.
  \item Se verifica que la lista de identificadores de clientes en el cl\'uster base se actualice adecuadamente.
  \item Se prueba que la fusi\'on no altere la integridad de los datos y que los puntos representativos del nuevo cl\'uster se calculen apropiadamente.
\end{itemize}

Casos bordes:

\begin{itemize}
  \item Manejo de listas vac\'ias de clientes en la creaci\'on.
  \item Fusi\'on de un cl\'uster consigo mismo (aunque no sea una operaci\'on v\'alida en la pr\'actica, se verifica que no lance excepciones inesperadas).
\end{itemize}

Estas pruebas aseguran que la estructura de datos subyacente al algoritmo CURE sean confiables y que las operaciones de agrupamiento se realicen de forma correcta. En la Figura~\ref{fig:prueba_merge} se puede apreciar el resultado proporcionado por el IDE a esta prueba.

\figura{fig17_prueba_merge.png}{0.85\textwidth}{Resultados de las pruebas de constructor y del m\'etodo merge.}{fig:prueba_merge}

\subsection{Pruebas del m\'etodo \textit{calculateClusterDistance} de la clase CURE}

El m\'etodo \textit{calculateClusterDistance} es uno de los m\'etodos m\'as importantes del algoritmo CURE, ya que determina la distancia entre dos cl\'usters utilizando los puntos representativos contra\'idos. Las pruebas unitarias dise\~nadas verifican:

\begin{itemize}
  \item C\'alculo de distancia entre cl\'usteres separados: se crean cl\'usteres con puntos en disposiciones horizontales y verticales, y se calcula la distancia esperada manualmente a partir de los puntos contra\'idos. Luego se compara con el resultado del m\'etodo, asegurando que coincida dentro de una tolerancia.
  \item Se prueba con cl\'usteres que tienen m\'ultiples puntos, comprobando que la distancia m\'inima entre los puntos contra\'idos sea la retornada.
  \item Simetr\'ia: se verifica que la distancia entre el cl\'uster A y el cl\'uster B sea igual a la distancia entre B y A.
  \item Manejo de cl\'usteres con un solo punto: en este caso, el centroide coincide con el punto, por lo que el punto contra\'ido es el mismo.
  \item Coordenadas negativas: se prueban cl\'usteres con puntos en coordenadas negativas para garantizar que el c\'alculo sea correcto en todo el plano.
  \item Uso de puntos representativos contra\'idos: se comprueba que los puntos utilizados en el c\'alculo sean efectivamente los puntos contra\'idos y no los originales.
  \item Valores positivos y finitos: se asegura que la distancia calculada sea siempre positiva y finita para cl\'usteres separados.
\end{itemize}

Estas pruebas garantizan que la m\'etrica de distancia entre cl\'usteres en el algoritmo CURE sea la adecuada para el proceso de agrupamiento jer\'arquico. En la Figura~\ref{fig:prueba_calcdist} se puede apreciar el resultado proporcionado por el IDE a esta prueba.

\figura{fig18_prueba_calcdist.png}{0.85\textwidth}{Resultados de las pruebas del m\'etodo calculateClusterDistance.}{fig:prueba_calcdist}

\section{Pruebas unitarias}

Adem\'as de las pruebas funcionales de los m\'etodos esenciales, se realizaron pruebas de caja negra para verificar el comportamiento global de los algoritmos DBSCAN y CURE en escenarios extremos, asegurando que las adaptaciones al MDVRP respetan las restricciones de capacidad y asignaci\'on.

En la Tabla~\ref{tab:pruebas_dbscan}, DBSCAN demuestra que su fase de expansi\'on controlada desde los dep\'ositos, junto con la reasignaci\'on final, logra respetar los l\'imites de capacidad incluso cuando la densidad de clientes supera la capacidad disponible. La Tabla~\ref{tab:pruebas_cure} muestra que CURE, mediante su mecanismo de fusi\'on jer\'arquica controlada, asigna los clientes de manera que no se exceda la capacidad de ning\'un dep\'osito, dejando como no asignados \'unicamente aquellos que no pueden ser absorbidos por ning\'un dep\'osito sin violar la restricci\'on. Estos resultados confirman que ambas adaptaciones cumplen con las reglas de negocio del MDVRP.

\begin{table}[H]
\centering
\caption{Pruebas unitarias de caja negra del algoritmo DBSCAN.}
\label{tab:pruebas_dbscan}
\footnotesize
\begin{tabular}{@{}cp{2.6cm}p{2.2cm}p{3cm}p{3cm}p{3cm}@{}}
\toprule
\textbf{No} & \textbf{Descripci\'on} & \textbf{Caso prueba} & \textbf{Juego de Datos} & \textbf{Resultados esperados} & \textbf{Resultados Reales} \\
\midrule
1 & Verificar asignaci\'on b\'asica de clientes a dep\'ositos cuando todos se encuentran dentro del radio $\varepsilon$. & Asignaci\'on b\'asica DBSCAN & 2 dep\'ositos y 6 clientes ubicados cerca de cada dep\'osito. Demandas peque\~nas dentro de la capacidad. & Cluster 7 = [1,2,3], Cluster 8 = [4,5,6], no asignados = 0 & Cluster 7 = [1,2,3], Cluster 8 = [4,5,6], no asignados = 0 \\
2 & Validar que el algoritmo respeta la capacidad m\'axima del dep\'osito. & Capacidad limitada DBSCAN & 1 dep\'osito con capacidad limitada y varios clientes cuya suma de demandas supera la capacidad. & Cluster 6 = [1,2], no asignados = [3,4,5] & Cluster 6 = [1,2], no asignados = [3,4,5] \\
3 & Validar la reasignaci\'on final de clientes no asignados. & Reasignaci\'on DBSCAN & Clientes alejados del dep\'osito inicial pero cercanos a otro dep\'osito con capacidad disponible. & Cluster 6 = [1,2,3], Cluster 7 = [4,5], no asignados = 0 & Cluster 6 = [1,2,3], Cluster 7 = [4,5], no asignados = 0 \\
4 & Validar excepci\'on cuando un deposito no puede alcanzar ning\'un cliente & Dep\'osito Vac\'io & Pocos clientes muy alejados de los dep\'ositos. & ClusterExcepci\'on & ClusterExcepci\'on \\
5 & Validar excepci\'on cuando los datos son err\'oneos. & Error en los Datos & Instancia sin dep\'ositos. & ProblemException & ProblemException \\
\bottomrule
\end{tabular}
\end{table}

\begin{table}[H]
\centering
\caption{Pruebas unitarias de caja negra del algoritmo CURE.}
\label{tab:pruebas_cure}
\footnotesize
\begin{tabular}{@{}cp{2.6cm}p{2.2cm}p{3cm}p{3cm}p{3cm}@{}}
\toprule
\textbf{No} & \textbf{Descripci\'on} & \textbf{Caso prueba} & \textbf{Juego de Datos} & \textbf{Resultados esperados} & \textbf{Resultados Reales} \\
\midrule
1 & Verificar la asignaci\'on b\'asica de clientes a dep\'ositos mediante fusi\'on de cl\'usteres. & Asignaci\'on b\'asica CURE & 2 dep\'ositos y varios clientes distribuidos en dos regiones cercanas a cada dep\'osito. & Cluster 7 = [1,2,3], Cluster 8 = [4,5,6], No asignados: 0 & Cluster 7 = [1,2,3], Cluster 8 = [4,6,5], No asignados: 0 \\
2 & Validar que la fusi\'on de cl\'usteres respeta la capacidad del dep\'osito. & Capacidad CURE & Clientes cercanos a un dep\'osito cuya demanda acumulada excede la capacidad disponible. & Cluster 6: [1,2], No asignados: [3,4,5] & Cluster 6: [1,2], No asignados: [3,4,5] \\
3 & Verificar el comportamiento cuando no es posible fusionar m\'as cl\'usteres. & No asignables CURE & Clientes restantes cuya asignaci\'on provocar\'ia exceder la capacidad de todos los dep\'ositos. & Cluster 5 = [1], Cluster 6 = [3], No asignados: [2,4] & Cluster 5 = [1], Cluster 6 = [3], No asignados: [2,4] \\
4 & Validar excepci\'on cuando los datos son err\'oneos. & Error en los Datos & Instancia sin dep\'ositos. & ProblemException & ProblemException \\
\bottomrule
\end{tabular}
\end{table}

\section{Comparaci\'on entre los algoritmos de agrupamiento}

Con el objetivo de analizar el comportamiento y la calidad de las asignaciones generadas por los algoritmos implementados en BHAVRP, se dise\~na un experimento enfocado en instancias utilizadas en trabajos anteriores [14], [15]. Este experimento busca comparar el rendimiento de cada algoritmo de asignaci\'on en condiciones controladas, evaluando los resultados mediante el conjunto de m\'etricas internas incluidas dentro de la biblioteca. El estudio experimental se lleva a cabo empleando las 33 instancias originalmente introducidas por [57] en su trabajo sobre el MDVRP. Estas instancias tienen amplio reconocimiento en el \'ambito cient\'ifico de la optimizaci\'on de rutas. La Tabla~\ref{tab:instancias} detalla, para cada instancia del problema, el n\'umero total de clientes, la cantidad de dep\'ositos, la flota de veh\'iculos asignada a cada dep\'osito y su capacidad. Cabe destacar que aquellas instancias que figuran con el mismo identificador en la columna `ID' presentan id\'enticas caracter\'isticas en cuanto a n\'umero de clientes, dep\'ositos, veh\'iculos por dep\'osito y capacidad de los veh\'iculos.

Para el experimento se usaron solo los algoritmos de agrupamiento de BHAVRP ya mencionados en la secci\'on 2.3. Al ser todos estos algoritmos deterministas, cada algoritmo se ejecuta una sola vez por instancia, utilizando como tipo de distancia aproximada la distancia euclidiana. Para analizar la calidad de las asignaciones generadas por cada algoritmo, se utilizan todas las m\'etricas de validaci\'on interna incorporadas en BHAVRP. Estas m\'etricas permiten evaluar cuantitativamente aspectos clave como la cohesi\'on y la separaci\'on entre cl\'usteres. En la Tabla~\ref{tab:metricas}, se presentan las m\'etricas a utilizar con sus siglas, el criterio que eval\'uan y una descripci\'on.

\begin{table}[H]
\centering
\caption{Descripci\'on de las instancias.}
\label{tab:instancias}
\small
\begin{tabular}{@{}ccccc@{}}
\toprule
\textbf{ID} & \textbf{\shortstack{Total de\\clientes}} & \textbf{\shortstack{Total de\\dep\'ositos}} & \textbf{\shortstack{Total de veh\'iculos\\por dep\'osito}} & \textbf{\shortstack{Capacidad\\de los veh\'iculos}} \\
\midrule
p01 & 50 & 4 & 4 & 80 \\
p02 & 50 & 4 & 2 & 160 \\
p03 & 75 & 5 & 3 & 140 \\
p04 & 100 & 2 & 8 & 100 \\
p05 & 100 & 2 & 5 & 200 \\
p06 & 100 & 3 & 6 & 100 \\
p07 & 100 & 4 & 4 & 100 \\
p08 & 249 & 2 & 14 & 500 \\
p09 & 249 & 3 & 12 & 500 \\
p10 & 249 & 4 & 8 & 500 \\
p11 & 249 & 5 & 6 & 500 \\
p12, p13, p14 & 80 & 2 & 5 & 60 \\
p15, p16, p17 & 160 & 4 & 5 & 60 \\
p18, p19, p20 & 240 & 6 & 5 & 60 \\
p21, p22, p23 & 360 & 9 & 5 & 60 \\
pr1 & 48 & 4 & 1 & 200 \\
pr2 & 96 & 4 & 2 & 195 \\
pr3 & 144 & 4 & 3 & 190 \\
pr4 & 192 & 4 & 4 & 185 \\
pr5 & 240 & 4 & 5 & 180 \\
pr6 & 288 & 4 & 6 & 175 \\
pr7 & 72 & 1 & 6 & 200 \\
pr8 & 144 & 2 & 6 & 190 \\
pr9 & 216 & 6 & 3 & 180 \\
pr10 & 288 & 6 & 4 & 170 \\
\bottomrule
\end{tabular}
\end{table}

\begin{table}[H]
\centering
\caption{Resumen de las m\'etricas incorporadas en BHAVRP.}
\label{tab:metricas}
\small
\begin{tabular}{@{}lp{1.3cm}p{2cm}p{8.5cm}@{}}
\toprule
\textbf{M\'etrica} & \textbf{Siglas} & \textbf{?`Qu\'e eval\'ua?} & \textbf{Descripci\'on} \\
\midrule
\textbf{Suma de Errores al Cuadrado} & SSE & Cohesi\'on & Suma de las distancias cuadr\'aticas entre cada punto y el centroide de su grupo. \textbf{Minimizar}. \\
\textbf{\'Indice de Ball-Hall} & BH & Cohesi\'on & Mide la dispersi\'on promedio dentro de los cl\'usteres. \textbf{Minimizar}. \\
\textbf{Suma de Cuadrados entre Cl\'usteres} & SSB & Separaci\'on & Eval\'ua la separaci\'on entre los grupos considerando la distancia entre sus centroides y el centro global de los datos. \textbf{Maximizar}. \\
\textbf{\'Indice de Calinski-Harabasz} & CH & Cohesi\'on y Separaci\'on & Relaci\'on entre la dispersi\'on entre cl\'usteres (SSB) y la dispersi\'on dentro de los cl\'usteres (SSE). \textbf{Maximizar}. \\
\textbf{\'Indice de Dunn} & DI & Cohesi\'on y Separaci\'on & Relaci\'on entre la distancia m\'inima inter-cl\'uster y el di\'ametro m\'aximo intra-cl\'uster (el tama\~no del grupo m\'as grande). \textbf{Maximizar}. \\
\textbf{\'Indice de Davies-Bouldin} & DB & Cohesi\'on y Separaci\'on & Analiza la relaci\'on entre la dispersi\'on interna de los cl\'usteres y la distancia entre sus centroides. \textbf{Minimizar}. \\
\textbf{Coeficiente de Silueta} & CS & Cohesi\'on y Separaci\'on & Combina cohesi\'on y separaci\'on para cada punto. \textbf{Maximizar}. \\
\textbf{\'Indice de Xie-Beni} & XB & Cohesi\'on y Separaci\'on & Penaliza los agrupamientos donde los puntos est\'an alejados de sus centroides y los que est\'an muy pr\'oximos entre s\'i. \textbf{Minimizar}. \\
\textbf{\'Indice de Ray-Turi} & RT & Cohesi\'on y Separaci\'on & Eval\'ua la relaci\'on entre la dispersi\'on dentro de los cl\'usteres y la distancia m\'inima entre los centroides de los cl\'usteres. \textbf{Minimizar}. \\
\bottomrule
\end{tabular}
\end{table}

\subsection{An\'alisis estad\'istico}

A partir de los experimentos realizados, se decide utilizar la prueba estad\'istica no param\'etrica de Friedman N $\times$ N [58], la cual posibilita enfrentar por pares todos los algoritmos examinados.

En la Tabla~\ref{tab:friedman} se presenta para cada algoritmo su valor promedio por cada m\'etrica en todas las instancias. Las celdas destacadas en verde reflejan las mejores posiciones, mientras que las celdas en rojo indican los desempe\~nos m\'as bajos. Al observar los rankings promedio se distingue una separaci\'on en el rendimiento de los algoritmos analizados. Este examen se ve complementado por los an\'alisis post-hoc.

\begin{table}[H]
\centering
\caption{Resultados de la prueba de Friedman NxN.}
\label{tab:friedman}
\footnotesize
\begin{tabular}{@{}lcccccccccc@{}}
\toprule
\textbf{Algoritmo} & \textbf{SSE} & \textbf{BH} & \textbf{SSB} & \textbf{CH} & \textbf{DI} & \textbf{CS} & \textbf{DB} & \textbf{XB} & \textbf{RT} & \textbf{Tiempo} \\
\midrule
\textbf{CURE} & \cellcolor{green!25}2.0303 & \cellcolor{green!25}2.2424 & \cellcolor{green!25}3.0909 & \cellcolor{green!25}2.1515 & 2.7879 & \cellcolor{green!25}2.2424 & \cellcolor{green!25}2.5152 & \cellcolor{green!25}2.5303 & \cellcolor{green!25}2.5303 & 2.0455 \\
\textbf{DBSCAN} & 4.4697 & 4.3939 & 3.9242 & 4.6818 & \cellcolor{red!20}5.4242 & 4.6212 & 4.2879 & 4.0455 & 4.0455 & \cellcolor{green!25}1.0455 \\
\textbf{Farthest-First} & 3.3333 & 3.3636 & 3.3939 & 3.1818 & 3.5758 & 3.1364 & 3.2424 & 3.1212 & 3.1212 & 4.1970 \\
\textbf{KMEANS} & 3.0606 & 3.1818 & 3.4242 & 2.9091 & 3.5606 & 2.8788 & 3.0909 & 2.9545 & 2.9545 & 3.8182 \\
\textbf{UPGMC} & 3.7121 & 3.3030 & 4.1364 & 3.9242 & \cellcolor{green!25}1.3030 & 3.8636 & 3.7727 & 4.0455 & 4.0455 & 3.9848 \\
\textbf{PAM} & 5.3636 & 5.3636 & 4.5152 & 5.1818 & 5.3636 & 5.1364 & 5.0303 & \cellcolor{red!20}5.3333 & \cellcolor{red!20}5.3333 & \cellcolor{red!20}7 \\
\textbf{CLARA} & \cellcolor{red!20}6.0303 & \cellcolor{red!20}6.1515 & \cellcolor{red!20}5.5152 & \cellcolor{red!20}5.9697 & \cellcolor{red!20}5.9848 & \cellcolor{red!20}6.1212 & \cellcolor{red!20}6.0606 & 5.9697 & 5.9697 & 5.9091 \\
\bottomrule
\end{tabular}
\end{table}

A continuaci\'on, se presenta la interpretaci\'on de los resultados obtenidos en el ranking:

\begin{itemize}
  \item CURE obtiene los mejores rankings en las m\'etricas de cohesi\'on y en la \'unica de separaci\'on obtuvo el mejor. Tambi\'en obtuvo el mejor ranking en la mayor\'ia de las combinadas, situ\'andose adem\'as como el segundo m\'as r\'apido. Esto indica que el nuevo algoritmo jer\'arquico es capaz de generar agrupaciones compactas y bien separadas en un buen tiempo.
  \item UPGMC presenta un comportamiento muy bueno en separaci\'on, pero sus rankings en cohesi\'on y m\'etricas combinadas son peores que los de CURE. Esto sugiere que, aunque UPGMC logra cl\'usteres muy separados, la consistencia interna de sus grupos es inferior.
  \item DBSCAN destaca por ser el algoritmo m\'as r\'apido. Sin embargo, muestra un desempe\~no moderado en cohesi\'on y pobre en separaci\'on y m\'etricas combinadas.
  \item K-Means y Farthest-First mantienen posiciones intermedias en todas las categor\'ias, mientras PAM y CLARA ocupan las \'ultimas posiciones, reflejando peor calidad de agrupamiento y adem\'as los peores tiempos.
\end{itemize}

Una vez obtenidos los rankings promedio de Friedman, se aplicaron los procedimientos \textit{post-hoc} de Bergmann [59], Holm [60] y Shaffer [61] para controlar la tasa de error familiar. A continuaci\'on, en la Tabla~\ref{tab:posthoc} se resumen los resultados de los procedimientos \textit{post hoc} de Bergmann, Holm y Shaffer para cada m\'etrica evaluada. Se indican las diferencias estad\'isticamente significativas ($\alpha=0.05$) y m\'as relevantes para el an\'alisis. Estos se emplearon porque son m\'etodos est\'andar para el control de rechazar una hip\'otesis nula verdadera tras un test de Friedman. Holm se usa por su simplicidad y aceptaci\'on general; Shaffer mejora la potencia al explotar las relaciones l\'ogicas entre las hip\'otesis de comparaciones por pares; y Bergmann proporciona la m\'axima potencia al determinar conjuntos de hip\'otesis aceptables.

\begin{table}[H]
\centering
\caption{Resumen an\'alisis post-hoc.}
\label{tab:posthoc}
\footnotesize
\begin{tabular}{@{}lp{4.5cm}lll@{}}
\toprule
\textbf{M\'etrica} & \textbf{Comparaci\'on} & \textbf{Holm} & \textbf{Shaffer} & \textbf{Bergmann} \\
\midrule
\textbf{SSE} & CURE $>$ CLARA, PAM & 0 & 0 & 0 \\
& CURE $>$ UPGMC & 0.020341 & 0.017212 & 0.014082 \\
& CURE $>$ DBSCAN & 0.000076 & 0.000067 & 0.000049 \\
& KMEANS $>$ CLARA & 0 & 0 & 0 \\
& KMEANS $>$ PAM & 0.000223 & 0.000223 & 0.000149 \\
& FF $>$ CLARA & 0.000007 & 0.000006 & 0.000004 \\
\textbf{DI} & UPGMC $>$ CURE & 0.036664 & 0.036664 & 0.026189 \\
& UPGMC $>$ CLARA, PAM & 0 & 0 & 0 \\
& CURE $>$ CLARA & 0 & 0 & 0 \\
& CURE $>$ PAM & 0.000020 & 0.000019 & 0.000009 \\
\textbf{CH} & CURE $>$ CLARA, PAM & 0 & 0 & 0 \\
& CURE $>$ UPGMC & 0.011155 & 0.009439 & 0.007723 \\
& KMEANS $>$ CLARA & 0 & 0 & 0 \\
& FF $>$ CLARA & 0.000003 & 0.000002 & 0.000002 \\
\textbf{CS} & CURE $>$ CLARA, PAM & 0 & 0 & 0 \\
& CURE $>$ UPGMC & 0.027603 & 0.025303 & 0.020703 \\
& KMEANS $>$ CLARA & 0 & 0 & 0 \\
\textbf{DB} & CURE, KMEANS $>$ CLARA & 0 & 0 & 0 \\
& FF $>$ CLARA & 0.000002 & 0.000002 & 0.000001 \\
\textbf{XB y RT} & CURE, KMEANS $>$ CLARA & 0 & 0 & 0 \\
& FF $>$ CLARA & 0.000002 & 0.000001 & 0.000001 \\
\textbf{Tiempo} & DBSCAN $>$ PAM, CLARA, FF & 0 & 0 & 0 \\
& DBSCAN $>$ UPGMC & 0.000001 & 0 & 0 \\
& DBSCAN $>$ KMEANS & 0.000002 & 0.000002 & 0.000001 \\
\bottomrule
\end{tabular}
\end{table}

El an\'alisis \textit{post-hocs} muestra las siguientes conclusiones:

\begin{itemize}
  \item CURE es el algoritmo con mejor calidad de agrupamiento, superando significativamente a UPGMC en la mayor\'ia de las m\'etricas combinadas y mostrando una clara ventaja en cohesi\'on y tiempo de ejecuci\'on.
  \item UPGMC solo conserva una ventaja significativa en el \'Indice de Dunn, lo que confirma que su buen comportamiento previo se debe a sus caracter\'isticas particulares y no a su naturaleza jer\'arquica.
  \item DBSCAN domina en tiempo de ejecuci\'on, pero su calidad es moderada y significativamente inferior a CURE en la mayor\'ia de los aspectos.
  \item KMEANS y Farthest-First ocupan posiciones intermedias, superando a PAM y CLARA.
  \item PAM y CLARA son consistentemente los peores en calidad y tiempo, siendo superados por todos los dem\'as algoritmos en m\'ultiples m\'etricas.
\end{itemize}

Es importante destacar que DBSCAN, en su adaptaci\'on con restricciones de capacidad, resulta extremadamente r\'apido, pero produce agrupaciones de calidad moderada en t\'erminos de separaci\'on. Esto puede ser aceptable en contextos donde la velocidad computacional sea prioritaria y se pueda tolerar una asignaci\'on menos refinada. No obstante, para instancias donde se requiere alta calidad en la creaci\'on de rutas posterior, otros algoritmos como CURE o K-Means ser\'ian m\'as adecuados.

Por \'ultimo, los rankings muestran que CURE supera a UPGMC en cohesi\'on y en la mayor\'ia de las m\'etricas combinadas, mientras que UPGMC solo le gana en el \'indice de Dunn. Esta diferencia sugiere que la estrategia de m\'ultiples puntos representativos y contracci\'on de CURE permite una mejor representaci\'on de la forma de los cl\'usteres. Esto tiene como resultado grupos m\'as homog\'eneos internamente sin sacrificar en exceso la separaci\'on. No obstante, para confirmar si estas diferencias son estad\'isticamente significativas, se requiere un an\'alisis espec\'ifico entre ambos algoritmos, el cual se aborda en el pr\'oximo ep\'igrafe mediante la prueba estad\'istica no param\'etrica de Wilcoxon.

\section{Comparaci\'on entre los algoritmos de agrupamiento jer\'arquico}

En este experimento se compara directamente UPGMC y CURE utilizando la prueba de Wilcoxon. Esto permitir\'a determinar si las diferencias observadas en los rankings de Friedman son estad\'isticamente significativas para cada m\'etrica. La prueba estad\'istica no param\'etrica de Wilcoxon se utiliza para dos muestras relacionadas, en este caso para dos algoritmos ejecutados en las mismas instancias [62].

Con el objetivo de fortalecer el an\'alisis experimental, se incorporan 11 instancias adicionales, alcanzando un total de 44 instancias para la evaluaci\'on de estos dos algoritmos. Estas instancias provienen de la investigaci\'on realizada en [63], y se detallan en la Tabla~\ref{tab:nuevas_instancias}.

\begin{table}[H]
\centering
\caption{Descripci\'on de las nuevas instancias.}
\label{tab:nuevas_instancias}
\small
\begin{tabular}{@{}ccccc@{}}
\toprule
\textbf{ID} & \textbf{\shortstack{Total de\\clientes}} & \textbf{\shortstack{Total de\\dep\'ositos}} & \textbf{\shortstack{Total de veh\'iculos\\por dep\'osito}} & \textbf{\shortstack{Capacidad de los\\veh\'iculos}} \\
\midrule
gj1 & 50 & 2 & 2 & 492 \\
gj2 & 75 & 3 & 3 & 415 \\
gj3 & 100 & 2 & 2 & 841 \\
gj4 & 150 & 4 & 4 & 651 \\
gj5 & 199 & 5 & 5 & 758 \\
gj6 & 75 & 2 & 2 & 669 \\
gj7 & 125 & 3 & 3 & 768 \\
gj8 & 249 & 4 & 4 & 1219 \\
gj9 & 60 & 2 & 2 & 579 \\
gj10 & 80 & 3 & 3 & 507 \\
gj11 & 120 & 5 & 5 & 1000 \\
\bottomrule
\end{tabular}
\end{table}

Ambos algoritmos se ejecutaron bajo las mismas condiciones experimentales descritas en la secci\'on 3.3, y los resultados para cada m\'etrica de validaci\'on interna fueron sometidos a la prueba de Wilcoxon.

A continuaci\'on, en la Tabla~\ref{tab:wilcoxon}, se presentan e interpretan los resultados de la prueba de Wilcoxon para cada una de las m\'etricas.

\begin{table}[H]
\centering
\caption{Resumen de los resultados de Wilcoxon.}
\label{tab:wilcoxon}
\small
\begin{tabular}{@{}lp{3.2cm}p{3.2cm}p{5.5cm}@{}}
\toprule
\textbf{M\'etrica} & \textbf{\shortstack[l]{p-value\\(CURE vs UPGMC)}} & \textbf{\shortstack[l]{p-value\\(UPGMC vs CURE)}} & \textbf{Interpretaci\'on ($\alpha = 0.05$)} \\
\midrule
CH & 0.000000000122 & $\geq$ 0.2 & CURE es significativamente mejor \\
DI & $\geq$ 0.2 & 0.0000616 & UPGMC es significativamente mejor \\
CS & 0.0000000164 & $\geq$ 0.2 & CURE es significativamente mejor \\
SSB & 0.00129 & $\geq$ 0.2 & CURE es significativamente mejor \\
BH & $\geq$ 0.2 & $\geq$ 0.2 & No hay diferencia significativa \\
DB & 0.06164 & $\geq$ 0.2 & No significativo ($p > 0.05$) aunque hay tendencia \\
RT & 0.000178 & $\geq$ 0.2 & CURE es significativamente mejor \\
SSE & 0.00000000913 & $\geq$ 0.2 & CURE es significativamente mejor \\
Tiempo & 0.000000000000114 & $\geq$ 0.2 & CURE es significativamente mejor \\
XB & 0.000153 & $\geq$ 0.2 & CURE es significativamente mejor \\
\bottomrule
\end{tabular}
\end{table}

La Tabla~\ref{tab:wilcoxon} ofrece una visi\'on concluyente de la comparaci\'on entre ambos algoritmos jer\'arquicos. CURE emerge como el algoritmo dominante en este estudio comparativo, superando a UPGMC de manera significativa en siete de las nueve m\'etricas de calidad de agrupamiento y en el tiempo de ejecuci\'on.

\begin{itemize}
  \item CURE demuestra ser superior en las m\'etricas de SSE y SSB. As\'i como en Calinski-Harabasz (CH), el Coeficiente de Silueta (CS) y el \'Indice de Xie-Beni (XB). Esto indica que, en promedio, la estructura de agrupamiento creada por CURE es de mayor calidad que la de UPGMC.
  \item El \'unico aspecto donde UPGMC supera significativamente a CURE es en el \'Indice de Dunn (DI).
  \item En el \'Indice de Ball-Hall (BH), ambos algoritmos se comportan de manera estad\'isticamente similar. Por otro lado, en el \'Indice de Davies-Bouldin (DB), la superioridad de CURE no es tan significativa.
  \item La clara ventaja en t\'erminos de tiempo de ejecuci\'on es para CURE. No solo produce mejores agrupaciones en la mayor\'ia de los casos, sino que lo hace significativamente m\'as r\'apido.
\end{itemize}

La evidencia demuestra que la naturaleza jer\'arquica por s\'i sola no garantiza superioridad. El algoritmo UPGMC, muestra el mejor desempe\~no en la m\'etrica del \'Indice de Dunn que es separaci\'on entre los cl\'usteres m\'as cercanos, es superado de manera estad\'isticamente significativa por CURE. Adem\'as, CURE logra estos mejores resultados con un costo computacional significativamente menor.

\section{Conclusiones parciales}

A partir del estudio experimental realizado en este cap\'itulo, se pueden extraer las siguientes conclusiones:

\begin{itemize}
  \item Las pruebas unitarias aplicadas a los nuevos algoritmos (DBSCAN y CURE) confirmaron el correcto funcionamiento de sus m\'etodos fundamentales, garantizando su integraci\'on fiable en la biblioteca BHAVRP.
  \item CURE se destaca como una buena alternativa, con los mejores rankings globales en la mayor\'ia de las m\'etricas, mientras que DBSCAN ofrece la opci\'on m\'as r\'apida, aunque de calidad m\'as limitada.
  \item El an\'alisis detallado mediante la prueba de Wilcoxon entre los dos algoritmos jer\'arquicos determina que la buena actuaci\'on previa de UPGMC se deb\'ia a sus caracter\'isticas particulares y no a su naturaleza jer\'arquica. CURE, que presenta una mec\'anica basada en puntos representativos, demuestra ser estad\'isticamente superior en la mayor\'ia de las m\'etricas de calidad y en tiempo de ejecuci\'on.
\end{itemize}

%==============================================================
%                      CONCLUSIONES
%==============================================================
\chapter*{Conclusiones}
\addcontentsline{toc}{chapter}{Conclusiones}
\markboth{Conclusiones}{}

La presente investigaci\'on ha permitido abordar el problema planteado, logrando incorporar nuevos algoritmos de agrupamiento en la Biblioteca de Heur\'isticas de Asignaci\'on para Problemas de Planificaci\'on de Rutas de Veh\'iculos con M\'ultiples Dep\'ositos (BHAVRP). A partir del trabajo realizado, se pueden extraer las siguientes conclusiones generales.

El estudio bibliogr\'afico desarrollado evidencia que la fase de asignaci\'on de clientes a dep\'ositos en el MDVRP puede ser abordada eficazmente mediante t\'ecnicas de agrupamiento de la miner\'ia de datos, aprovechando su capacidad para descubrir estructuras en los datos geogr\'aficos. Asimismo, el an\'alisis comparativo de bibliotecas existentes reafirma que BHAVRP constituye una buena alternativa dentro del MDVRP, al ser la \'unica especializada en esta fase y con una arquitectura abierta y extensible que facilita la incorporaci\'on de nuevos algoritmos. La implementaci\'on de los algoritmos DBSCAN y CURE aumenta la diversidad de enfoques de la biblioteca. La incorporaci\'on de restricciones de capacidad y adaptar los mecanismos de expansi\'on y fusi\'on, demuestra que es posible aplicar t\'ecnicas cl\'asicas de miner\'ia de datos a contextos de optimizaci\'on log\'istica, garantizando la factibilidad de las soluciones generadas.

Por \'ultimo, el an\'alisis experimental desarrollado valida el correcto funcionamiento de los algoritmos implementados a trav\'es de pruebas unitarias. El primer experimento posiciona a CURE como el algoritmo de rendimiento superior, obteniendo los mejores rankings en la mayor\'ia de las m\'etricas de calidad de agrupamiento y un tiempo de ejecuci\'on aceptable. Por otra parte, DBSCAN demuestra ser la opci\'on m\'as r\'apida, aunque con una calidad de agrupamiento menor. En cuanto al segundo experimento, la comparaci\'on entre los dos algoritmos jer\'arquicos mediante la prueba de Wilcoxon demuestra que existen diferencias significativas. Se determina que la superioridad previamente observada en UPGMC no se debe a su naturaleza jer\'arquica, sino a sus caracter\'isticas particulares. CURE, con su estrategia de m\'ultiples puntos representativos y contracci\'on, resulta estad\'isticamente superior en la mayor\'ia de las m\'etricas de calidad y en tiempo de ejecuci\'on.

%==============================================================
%                      RECOMENDACIONES
%==============================================================
\chapter*{Recomendaciones}
\addcontentsline{toc}{chapter}{Recomendaciones}
\markboth{Recomendaciones}{}

A partir de los resultados alcanzados durante el desarrollo de esta investigaci\'on, se formulan las siguientes recomendaciones para trabajos futuros.

\begin{itemize}
  \item Se recomienda explorar la incorporaci\'on de otros algoritmos de agrupamiento basados en densidad, como OPTICS, que podr\'ia ofrecer una mejor respuesta ante la presencia de cl\'usteres con densidades variables, superando una de las limitaciones conocidas de DBSCAN.
  \item Dado que el presente estudio se ha centrado en la calidad de la asignaci\'on mediante m\'etricas de validaci\'on interna, se recomienda experimentar a trav\'es de herramientas como BHCVRP con estos dos nuevos algoritmos para ver la calidad de las rutas obtenidas y si se corresponden con lo prometedoras que son las asignaciones, CURE por calidad de las m\'etricas y DBSCAN por su rapidez.
  \item Por \'ultimo, se recomienda incorporar estos cambios a la versi\'on en Python de la biblioteca.
\end{itemize}

%==============================================================
%                 REFERENCIAS BIBLIOGRAFICAS
%==============================================================
\addcontentsline{toc}{chapter}{Referencias bibliogr\'aficas}
\begin{thebibliography}{99}
\markboth{Referencias bibliogr\'aficas}{}

\bibitem{1} M.~E. Coloma Salazar, J.~Arzola Ruiz, y C.~E. Marrero Fornaris, ``La gesti\'on de la distribuci\'on y la planificaci\'on de rutas en el sector empresarial cubano'', \emph{Rev. Univ. Soc.}, vol.~14, pp. 357--368, 2022.

\bibitem{2} G.~B. Dantzig y J.~H. Ramser, \emph{The Truck Dispatching Problem. Management Science}, 6.\textsuperscript{a} ed., 1959.

\bibitem{3} Clarke, G. y Wright, J.R., \emph{Scheduling of Vehicle Routing Problem from a Central Depot to a Number of Delivery Points. Operations Research}, vol.~12, 1964.

\bibitem{4} S.~Arora y B.~Barak, \emph{Computational complexity: a modern approach}. Cambridge University Press, 2009.

\bibitem{5} K.~Braekers, K.~Ramaekers, y I.~Van Nieuwenhuyse, ``The vehicle routing problem: State of the art classification and review'', \emph{Comput. Ind. Eng.}, vol.~99, pp. 300--313, sep. 2016, doi: 10.1016/j.cie.2015.12.007.

\bibitem{6} D.~I.~G. Mateos, ``Desarrollo de un m\'etodo h\'ibrido para la resoluci\'on del MDVRP'', \emph{Rev. Esc. Jacobea Posgrado}, 2013.

\bibitem{7} J.~A.~G. Bermudez, ``Aplicaci\'on de t\'ecnicas de clusterizaci\'on como inicializadores y la metaheur\'istica ILS, al Problema de Ruteo Veh\'iculos con M\'ultiples Dep\'ositos (MDVRP)'', Proyecto de grado presentado como requisito parcial para optar al t\'itulo de Magister en Ingenier\'ia El\'ectrica, Universidad Tecnol\'ogica de Pereira.

\bibitem{8} S.~Barreto, C.~Ferreira, J.~Paix\~ao, y B.~S. Santos, ``Using clustering analysis in a capacitated location-routing problem'', \emph{Eur. J. Oper. Res.}, vol.~179, n.\textsuperscript{o}~3, pp. 968--977, 2007, doi: 10.1016/j.ejor.2005.06.074.

\bibitem{9} Y.~Z. Mehrjerdi y A.~Nadizadeh, ``Using greedy clustering method to solve capacitated location-routing problem with fuzzy demands'', \emph{Eur. J. Oper. Res.}, vol.~229, n.\textsuperscript{o}~1, pp. 75--84, 2013, doi: 10.1016/j.ejor.2013.02.013.

\bibitem{10} A.~K. Jain y R.~C. Dubes, \emph{Algorithms for clustering data}. Prentice Hall advanced reference series. Englewood Cliffs, NJ: Prentice-Hall, 1988.

\bibitem{11} K.~Khan, S.~U. Rehman, K.~Aziz, S.~Fong, y S.~Sarasvady, ``DBSCAN: Past, present and future'', en \emph{The fifth international conference on the applications of digital information and web technologies (ICADIWT 2014)}, IEEE, 2014, pp. 232--238.

\bibitem{12} S.~Guha, R.~Rastogi, y K.~Shim, ``Cure: an efficient clustering algorithm for large databases'', \emph{Inf. Syst.}, vol.~26, n.\textsuperscript{o}~1, pp. 35--58, mar. 2001, doi: 10.1016/S0306-4379(01)00008-4.

\bibitem{13} ``GitHub - vrpsolutions/BHAVRP: BHAVRP: Componente de software para la asignaci\'on de clientes a dep\'ositos en problemas de rutas de veh\'iculos''. Accedido: 21 de enero de 2026. [En l\'inea]. Disponible en: https://github.com/vrpsolutions/BHAVRP

\bibitem{14} E.~R. Arag\'on, I.~T. P\'erez, A.~R. Su\'arez, y A.~de la C.~M. Morales, ``BHAVRP: Biblioteca de Heur\'isticas de Asignaci\'on para el Problema de Planificaci\'on de Rutas de Veh\'iculos con M\'ultiples Dep\'ositos'', \emph{Rev. Cuba. Ing. RCI}, pp. 1--13, 2025.

\bibitem{15} Eric Ramos Arag\'on, ``Incorporaci\'on de un m\'odulo de m\'etricas de evaluaci\'on para algoritmos de asignaci\'on en BHAVRP'', Trabajo para optar por el t\'itulo de Ingeniero Inform\'atico, Universidad Tecnol\'ogica de la Habana ``Jos\'e Antonio Echeverr\'ia'', La Habana, Cuba, 2025.

\bibitem{16} S.~N. Kumar, ``A survey on the vehicle routing problem and its variants'', \emph{Intell. Inf. Manag.}, 2012.

\bibitem{17} G.~Berbeglia, J.-F. Cordeau, y G.~Laporte, ``Dynamic pickup and delivery problems'', \emph{Eur. J. Oper. Res.}, vol.~202, n.\textsuperscript{o}~1, pp. 8--15, 2007, doi: 10.1016/j.ejor.2009.04.024.

\bibitem{18} R.~Baldacci, M.~Battarra, y D.~Vigo, ``Routing a heterogeneous fleet of vehicles'', en \emph{The Vehicle Routing Problem: Latest Advances and New Challenges}, B.~Golden, S.~Raghavan, y E.~Wasil, Eds., Springer, 2008, pp. 3--27. doi: 10.1007/978-0-387-77778-8\_1.

\bibitem{19} D.~Sariklis y S.~Powell, ``A heuristic method for the open vehicle routing problem'', \emph{J. Oper. Res. Soc.}, vol.~51, n.\textsuperscript{o}~5, pp. 564--573, 2000, doi: 10.1057/palgrave.jors.2600924.

\bibitem{20} J.-F. Cordeau, G.~Desaulniers, J.~Desrosiers, M.~M. Solomon, y F.~Soumis, ``The VRP with time windows'', en \emph{The Vehicle Routing Problem}, P.~Toth y D.~Vigo, Eds., SIAM, 2002, pp. 157--193. doi: 10.1137/1.9780898718515.ch7.

\bibitem{21} J.~R. Montoya-Torres, J.~L\'opez Franco, S.~Nieto Isaza, H.~Felizzola Jim\'enez, y N.~Herazo-Padilla, ``A literature review on the vehicle routing problem with multiple depots'', \emph{Comput. Ind. Eng.}, vol.~79, pp. 115--129, 2015, doi: 10.1016/j.cie.2014.10.029.

\bibitem{22} M.~Gendreau, G.~Laporte, y R.~S\'eguin, ``Stochastic vehicle routing'', \emph{Eur. J. Oper. Res.}, vol.~88, n.\textsuperscript{o}~1, pp. 3--12, 1996, doi: 10.1016/0377-2217(95)00050-X.

\bibitem{23} P.~M. Francis, K.~R. Smilowitz, y M.~Tzur, ``The period vehicle routing problem and its extensions'', en \emph{The Vehicle Routing Problem: Latest Advances and New Challenges}, B.~Golden, S.~Raghavan, y E.~Wasil, Eds., Springer, 2008, pp. 73--102. doi: 10.1007/978-0-387-77778-8\_4.

\bibitem{24} A.~Ibrahim, R.~Abdulaziz, y J.~Ishaya, ``Capacitated vehicle routing problem'', \emph{Int. J. Res. - GRANTHAALAYAH}, vol.~7, pp. 310--327, abr. 2019, doi: 10.5281/zenodo.2636820.

\bibitem{25} V.~Rebillas Loredo, ``The multi-depot VRP with vehicle interchanges'', PhD Thesis, Universitat Polit\`ecnica de Catalunya, 2018. doi: 10.5821/dissertation-2117-127490.

\bibitem{26} J.~B. Colin, ``Dise\~no de un algoritmo metaheur\'istico para la soluci\'on del Problema de Ruteo de Veh\'iculos MultiDep\'osito'', para obtener el grado de Maestr\'ia en Optimizaci\'on y C\'omputo Aplicado, Universidad Aut\'onoma del Estado de Morelos, 2024.

\bibitem{27} NEO Research Group, ``Multiple Depot VRP''. Universidad de M\'alaga, 2019. [En l\'inea]. Disponible en: https://neo.lcc.uma.es/vrp/vrp-flavors/multiple-depot-vrp/

\bibitem{28} T.~R.~P. Ramos, M.~I. Gomes, y A.~P. Barbosa-P\'ovoa, ``Multi-depot vehicle routing problem: a comparative study of alternative formulations'', \emph{Int. J. Logist. Res. Appl.}, vol.~23, n.\textsuperscript{o}~2, pp. 103--120, 2020, doi: 10.1080/13675567.2019.1630374.

\bibitem{29} F.~Hamid, Ed., \emph{Optimization essentials: theory, tools, and applications}. International series in operations research \& management science, vol.~353. Singapore: Springer, 2024.

\bibitem{30} M.~J\"unger \emph{et al.}, ``50 Years of Integer Programming 1958-2008: From the Early Years to the State-of-the-Art'', 2009, doi: 10.1007/978-3-540-68279-0.

\bibitem{31} D.~G. Mogale, A.~Ghadge, y S.~K. Jena, ``Modelling and optimising a multi-depot vehicle routing problem for freight distribution in a retail logistics network'', \emph{Comput. Ind. Eng.}, vol.~207, p.~111315, sep. 2025, doi: 10.1016/j.cie.2025.111315.

\bibitem{32} E.~Erdem, T.~Aydin, y B.~Erkayman, ``AI-Infused Evolutionary Optimization for Multi-Depot Vehicle Routing'', \emph{Case Study Beverage Distrib.}, vol.~37, n.\textsuperscript{o}~27-28, dic. 2025, doi: 10.1002/cpe.70431.

\bibitem{33} F.~S. Hern\'andez, ``Aplicaci\'on del modelo VRP (Vehicle Routing Problem) para la optimizaci\'on de una red de distribuci\'on'', Universidad Nacional Aut\'onoma de M\'exico, 2015.

\bibitem{34} I.-M. Chao, B.~L. Golden, y E.~A. Wasil, ``A new heuristic for the multi-depot vehicle routing problem that improves upon best-known solutions'', \emph{Am. J. Math. Manag. Sci.}, vol.~13, pp. 371--406, 1993.

\bibitem{35} F.~A. Tillman, ``The multiple terminal delivery problem with probabilistic demands'', \emph{Transp. Sci.}, vol.~3, n.\textsuperscript{o}~3, pp. 192--204, 1969, doi: 10.1287/trsc.3.3.192.

\bibitem{36} B.~E. Gillett y L.~R. Miller, ``A heuristic algorithm for the vehicle-dispatch problem'', \emph{Oper. Res.}, vol.~22, n.\textsuperscript{o}~2, pp. 340--349, 1974, doi: 10.1287/opre.22.2.340.

\bibitem{37} Y.~Shi, L.~Lv, F.~Hu, y Q.~Han, ``A Heuristic Solution Method for Multi-Depot Vehicle Routing-Based Waste Collection Problems'', \emph{Appl. Sci.}, vol.~10, n.\textsuperscript{o}~7, 2020, doi: 10.3390/app10072403.

\bibitem{38} C.~Blum y A.~Roli, ``Metaheuristics in combinatorial optimization: Overview and conceptual comparison'', \emph{ACM Comput. Surv.}, vol.~35, n.\textsuperscript{o}~3, pp. 268--308, sep. 2003, doi: 10.1145/937503.937505.

\bibitem{39} J.~H. Orallo, M.~J.~R. Quintana, y C.~F. Ram\'irez, \emph{Introducci\'on a la Miner\'ia de Datos}. Pearson Educaci\'on, 2004.

\bibitem{40} G.~Kesavaraj y S.~Sukumaran, ``A study on classification techniques in data mining'', en \emph{2013 fourth international conference on computing, communications and networking technologies (ICCCNT)}, IEEE, 2013, pp. 1--7.

\bibitem{41} \c{S}. \c{S}eker, A.~F. G\"uneri, G.~Ermi\c{s}, y \.{I}. Y{\i}ld{\i}z, ``Solving the multi depot vehicle routing problem with assignment algorithm'', en \emph{13th International Research/Expert Conference ``Trends in the Development of Machinery and Associated Technology'', Hammamet}, 2009, pp. 225--228.

\bibitem{42} K.~Koutroumbas y S.~Theodoridis, \emph{Pattern recognition}. Academic Press, 2008.

\bibitem{43} S.~Salhi y M.~Sari, ``A heuristic algorithm for the multi-depot vehicle routing problem'', \emph{J. Oper. Res. Soc.}, vol.~49, n.\textsuperscript{o}~4, pp. 424--431, 1998, doi: 10.1057/palgrave.jors.2600526.

\bibitem{44} JSprit Contributors, ``JSprit: A Java-based library for solving vehicle routing problems''. Accedido: 18 de marzo de 2026. [En l\'inea]. Disponible en: https://jsprit.readthedocs.io/en/latest/

\bibitem{45} L.~Perron y V.~Furnon, \emph{OR-Tools}. Google. [En l\'inea]. Disponible en: https://developers.google.com/optimization/

\bibitem{46} F.~Pedregosa, \emph{Scikit-learn: Machine Learning in Python}, 12.\textsuperscript{a} ed., 2011.

\bibitem{47} Ananda de la Caridad Morales Morales, ``Nueva versi\'on de la biblioteca de heur\'isticas de construcci\'on para problemas de planificaci\'on de rutas de veh\'iculos'', Trabajo para optar por el t\'itulo de Ingeniero Inform\'atico, Universidad Tecnol\'ogica de la Habana ``Jos\'e Antonio Echeverr\'ia'', La Habana, Cuba, 2025.

\bibitem{48} G.~Ya\~nez y I.~Torres, ``Desarrollo de nuevos modelos para el Problema del Transporte Obrero en CaSVRP'', Universidad Tecnol\'ogica de la Habana ``Jos\'e Antonio Echeverr\'ia'', La Habana, Cuba, 2021.

\bibitem{49} R.~D. Hern\'andez, ``Nueva versi\'on de la biblioteca de clases BiCIAM para solucionar problemas multiobjetivo'', Trabajo de diploma, Universidad Tecnol\'ogica de La Habana ``Jos\'e Antonio Echeverr\'ia'', La Habana, Cuba, 2014.

\bibitem{50} R.~D. Hern\'andez, ``Nueva versi\'on de la biblioteca de clases BiCIAM para solucionar problemas multiobjetivo'', Trabajo de diploma (Ingenier\'ia Inform\'atica), Universidad Tecnol\'ogica de La Habana ``Jos\'e Antonio Echeverr\'ia'', La Habana, Cuba, 2014.

\bibitem{51} L.~Lores, ``Operadores de mutaci\'on basados en heur\'isticas de construcci\'on para Problemas de Planificaci\'on de Rutas de Veh\'iculos'', Trabajo de diploma (Ingenier\'ia Inform\'atica), Universidad Tecnol\'ogica de La Habana ``Jos\'e Antonio Echeverr\'ia'', La Habana, Cuba, 2017.

\bibitem{52} W.~Quiala Fonseca, ``Agrupamiento de datos desde un enfoque paralelo'', \emph{Rev. Cuba. Cienc. Inform\'aticas}, vol.~17, 2023.

\bibitem{53} A.~Patra, R.~Khan, y S.~Vijayalakshmi, ``Customer segmentation and future purchase prediction using RFM measures'', en \emph{2022 4th International Conference on Advances in Computing, Communication Control and Networking (ICAC3N)}, IEEE, 2022, pp. 753--759.

\bibitem{54} M.~Ester, H.-P. Kriegel, J.~Sander, X.~Xu, y others, ``A density-based algorithm for discovering clusters in large spatial databases with noise'', en \emph{kdd}, 1996, pp. 226--231.

\bibitem{55} S.~Suman, A.~Karna, y K.~Gibert, ``Bootstrap-CURE: A Novel Clustering Approach for Sensor Data---An Application to 3D Printing Industry'', \emph{Appl. Sci.}, vol.~12, n.\textsuperscript{o}~4, 2022, doi: 10.3390/app12042191.

\bibitem{56} W.~Wu, H.~Xiong, y S.~Shekhar, \emph{Clustering and information retrieval}, vol.~11. Springer Science \& Business Media, 2003.

\bibitem{57} J.~F. Cordeau, J.~Desrosiers, y M.~M. Solomon, ``The VRP with Multiple Depots: State-of-the-art and New Insights'', \emph{Transportation Science}, vol.~36, n.\textsuperscript{o}~4, pp. 278--292, 2002.

\bibitem{58} M.~Friedman, ``The use of ranks to avoid the assumption of normality implicit in the analysis of variance'', \emph{J. Am. Stat. Assoc.}, vol.~32, n.\textsuperscript{o}~200, pp. 675--701, 1937.

\bibitem{59} B.~Bergmann y G.~Hommel, ``Improvements of General Multiple Test Procedures for Redundant Systems of Hypotheses'', en \emph{Multiple Hypothesenpr\"ufung --- Multiple Hypotheses Testing}, Berlin, Heidelberg: Springer-Verlag, 1988, pp. 100--115. doi: 10.1007/978-3-642-52307-6\_8.

\bibitem{60} S.~Holm, ``A Simple Sequentially Rejective Multiple Test Procedure'', \emph{Scand. J. Stat.}, vol.~6, pp. 65--70, 1979.

\bibitem{61} J.~P. Shaffer, ``Modified Sequentially Rejective Multiple Test Procedures'', \emph{J. Am. Stat. Assoc.}, vol.~81, n.\textsuperscript{o}~395, pp. 826--831, 1986, doi: 10.1080/01621459.1986.10478341.

\bibitem{62} F.~Wilcoxon, ``Individual Comparisons by Ranking Methods'', \emph{Biom. Bull.}, vol.~1, n.\textsuperscript{o}~6, pp. 80--83, 1945, doi: 10.2307/3001968.

\bibitem{63} G.~E. Gillett y L.~R. Johnson, ``Multi-terminal vehicle-dispatch algorithm'', \emph{Transp. Res.}, vol.~10, n.\textsuperscript{o}~1, pp. 47--55, 1976.

\end{thebibliography}

\end{document}